\documentclass[11pt,a4paper]{article}
\usepackage[margin=25mm]{geometry}
\usepackage[T1]{fontenc}
\usepackage{lmodern,microtype}
\usepackage{amsmath,amssymb,amsthm,mathtools}
\usepackage{xcolor}
\usepackage[colorlinks=true,linkcolor=blue!45!black,citecolor=blue!45!black,urlcolor=blue!45!black]{hyperref}
\usepackage{fancyhdr}
\setlength{\headheight}{14pt}
\pagestyle{fancy}\fancyhf{}
\fancyhead[L]{\small A quartic construction in dimension two}
\fancyhead[R]{\small Research draft}
\fancyfoot[C]{\thepage}
\setlength{\parskip}{2pt}
\setlength{\parindent}{14pt}
\newtheorem{theorem}{Theorem}[section]
\newtheorem{lemma}[theorem]{Lemma}
\newtheorem{proposition}[theorem]{Proposition}
\newtheorem{corollary}[theorem]{Corollary}
\theoremstyle{remark}
\newtheorem{remark}[theorem]{Remark}
\newcommand{\Q}{\mathbf Q}
\newcommand{\Z}{\mathbf Z}
\newcommand{\F}{\mathbf F}
\newcommand{\A}{\mathbf A}
\newcommand{\PP}{\mathbf P}
\newcommand{\OO}{\mathcal O}
\newcommand{\Spec}{\operatorname{Spec}}
\newcommand{\Frac}{\operatorname{Frac}}
\newcommand{\Pic}{\operatorname{Pic}}
\newcommand{\Cl}{\operatorname{Cl}}
\newcommand{\colim}{\operatorname*{colim}}
\begin{document}
\thispagestyle{plain}
\begin{center}
{\LARGE A quartic construction for rational\\[3pt]
Gersten injectivity in dimension two}\\[8pt]
{\normalsize Actual arithmetic classes and a continuous elliptic detector}\\[9pt]
{\small Proposed counterexample manuscript --- 7 October 2026}
\end{center}
\vspace{4pt}
\noindent\fcolorbox{black!35}{black!3}{%
\parbox{\dimexpr\textwidth-2\fboxsep-2\fboxrule\relax}{\small
\textbf{Review status.} This manuscript presents a new proposed
counterexample, developed and subjected to internal adversarial
checking during the current research session. It has not been
externally refereed or formally verified. The references establish
the stated inputs; they do not establish the new conclusion.
The separate three-dimensional constructions do not depend on
this dimension-two argument.}}

\begin{abstract}
We give a proposed construction of an infinite-order class in the
kernel of the rational generic map on $K_3$ for a two-dimensional
ramified regular local ring, and for its maximal-ideal completion.
The residue field is $\mathbf F_5(T)$.  An actual cyclotomic
$K_3$-coefficient is pushed forward from a regular principal divisor.
Its image after a scaled ramified base change is detected by a
weighted root-divisor class.  The corresponding quartic has a
constant ordinary elliptic quotient.  A bounded contraction in
the elliptic direction and a scalar Frobenius pole argument show
that the Chern class survives completion at the generic Cohen
point.  The relative Beilinson fiber comparison, together with two
separate localization injections, turns this computation into
nonvanishing of the actual arithmetic class. A second residue
parameter and a strict Laurent-residue map yield one fixed complete
regular local ring of dimension two whose rational generic kernels
are infinite-dimensional in every degree at least three.
\end{abstract}

\paragraph{Related recent preprints.}
The preprints \cite{PriorIntegralK3,PriorIntegralK5} claim integral
counterexamples in degrees three and five on different ramified
rings with finite residue fields. Both explicitly leave the order
of their constructed class, and its image after rationalization,
undetermined. The present argument requires infinite order and
uses a transcendental residue parameter in dimension two. Those
preprints are not inputs to the proof; their cited claims are not
treated here as externally verified. No priority claim is made.


\section{The proposed statement and the actual class}
\label{sec:dimtwoStatement}

All $K$-groups are algebraic $K$-groups.  For an abelian group $M$,
write $M_{\Q}=M\otimes_{\Z}\Q$.  The cyclotomic polynomial is
\[
 \Phi_5(1+y)=5+10y+10y^2+5y^3+y^4.
\]
The variable $T$ is a coefficient parameter that remains
transcendental in the residue field.

\begin{theorem}[Proposed dimension-two counterexample]
\label{thm:dimensiontwoMain}
Put
\[
 A_2=\left(\Z_{(5)}[T,x,y]/
   (\Phi_5(1+y)+x^4+Txy^3)\right)_{(5,x,y)},
 \qquad \widehat A_2=\widehat{A_2}^{\,(x,y)}.
\]
Both are regular local domains of dimension two, with maximal
ideal $(x,y)$ and residue field $\F_5(T)$.  There is an actual
class $c\in K_3(A_2)$ such that its image
$\widehat c\in K_3(\widehat A_2)$ has infinite order, and
\[
 c|_{\Frac A_2}=0,\qquad
 \widehat c|_{\Frac\widehat A_2}=0
\]
integrally.  In particular, both rational generic maps in degree
three have nonzero kernel.
\end{theorem}

Here is the supported class used in the theorem.  Let $\rho$ be
a primitive fifth root of unity and set
$W_{\rm ar}=\Z_{(5)}[\rho]$, $\pi=\rho-1$.
The divisor $x=0$ has coordinate ring
\[
 D_{\rm ar}=A_2/(x)\simeq W_{\rm ar}[T]_{(\pi)},
 \qquad 1+y\longmapsto\rho.
\]
It is a regular DVR.  Start with the rational weight-two
root-of-unity class $b_2(-\rho)$, lift it to the coefficient DVR,
and clear one integer denominator to get an actual
$\beta_0\in K_3(W_{\rm ar})$.  Define
\[
 \beta=(1-\sigma_{-1})\beta_0,
 \qquad \sigma_{-1}(\rho)=\rho^{-1},
 \qquad c=i_*(\beta|_{D_{\rm ar}}),
\]
where $i:\Spec D_{\rm ar}\hookrightarrow\Spec A_2$.
The pushforward is defined in $G$-theory and identified with
$K$-theory by regularity.  Its generic vanishing follows directly
from localization at $x$; the work lies in proving infinite order.

Section~\ref{sec:dimtwoExplicitTriple} gives a second, completely
explicit representative of a nonzero rational kernel class:
\[
 c_{\rm exp}=\left\{2+y,\ 1-x,\
 1+(1-x)(x^3+Ty^3)y\bigl((1+y)^5+1\bigr)\right\}.
\]
The braces denote the product of three unit classes in algebraic
$K$-theory. Its vanishing after inverting $x$ is an integral
Dennis--Stein identity. Its infinite order follows by identifying
its image on the specified test algebra with a nonzero rational
multiple of the detected weighted class.

Section~\ref{sec:dimtwoVisibleHigherRank} extends the detector to
all degrees $K_{2n-1}$, $n\ge2$. On one fixed complete regular local
ring of dimension two and residue field $\overline{\F}_5(T)$, it
produces infinite-dimensional rational generic kernels in every
such degree. Section~\ref{sec:dimtwoAllDegreesFixedRing} then adds
a second residue parameter and constructs a single complete regular
local ring with residue field $\overline{\F}_5(T,S)$ whose rational
generic kernels are infinite-dimensional in every degree $K_j$,
$j\ge3$. The Laurent residue detects the even products; their
independence also proves independence of the odd factors.

The proof is organized around three specific issues.  First,
the two odd character components of this \emph{actual}
coefficient are nonzero.  Second, the root Chern class must
remain nonzero after completing an infinite localization of the
parameter line; a scalar Frobenius equation controls possible
new primitives.  Third, a nonzero relative completed character
must be transported back to the actual source class.  The last
step uses vanishing in two adjacent degrees of the special-fiber
$K$-spectrum, and does not assume that actual $K$-theory maps
injectively to its completion.

The finite-type source and its completion are handled in the
same argument.  In the completed differential complexes,
cohomology always uses the ordinary image of the differential,
without taking its closure.  This distinction is essential for
the unit-root class that appears below.



% Root-written arithmetic section for the separate dimension-two research draft.
% This is a new proposed argument, not an edit of the three-dimensional paper.
\section{A cyclotomic coefficient in both odd characters}
\label{sec:dimtwoExplicitCoefficient}

Put
\[
 W_{\rm ar}=\mathbf Z_{(5)}[\rho],\qquad
 W_E=\mathbf Z_5[\rho],\qquad E=\mathbf Q_5(\rho),\qquad
 \rho^5=1,\quad \rho\ne1,
\]
and write $\pi=\rho-1$.  Both coefficient rings are DVRs, with
uniformizer $\pi$ and residue field $\mathbf F_5$.  Let
$\sigma_g(\rho)=\rho^g$, for $g\in\mathbf F_5^\times$, and let
$\omega:\mathbf F_5^\times\to\mathbf Z_5^\times$ be the Teichmuller
character.  For $r=1,3$ define operators only on the regulator target:
\[
 e_r=\frac14\sum_{g=1}^4\omega(g)^{-r}\sigma_g,
 \qquad \tau_r=\sum_{g=1}^4\omega(g)^{-r}\rho^g.
\]
No $\mathbf Q_5$-linear projector will be applied inside an actual
rational $K$-group.

\begin{lemma}[Two explicit nonzero character components]
\label{lem:dimtwoDilogCharacters}
For the Coleman dilogarithm normalized by its series at zero,
\[
 \frac{4e_1\operatorname{Li}_2(-\rho)}{\tau_1}
       \in 2+5\mathbf Z_5,
 \qquad
 \frac{4e_3\operatorname{Li}_2(-\rho)}{\tau_3}
       \in 1+5\mathbf Z_5.
\]
Moreover $v_\pi(\tau_r)=r$.  In particular both projected values
are nonzero.
\end{lemma}

\begin{proof}
Define the rational functions
\[
 S_j(z)=\left(z\frac{d}{dz}\right)^j\frac1{1-z}
       \in\mathbf Z[z,(1-z)^{-1}].
\]
On $|z|<1$ they have the series
$S_j(z)=\sum_{m\ge0}m^jz^m$, with $0^0=1$ for $j=0$.
Their values at $-1$ below are values of these rational functions.
The depleted dilogarithm satisfies
\begin{align}
 F(z)&=\operatorname{Li}_2(z)-5^{-2}\operatorname{Li}_2(z^5)
 \label{eq:dimtwoDepletion}\\
 &=\sum_{j\ge0}(-1)^j(j+1)5^j S_j(z^5)
                  \sum_{a=1}^4\frac{z^a}{a^{j+2}}.\notag
\end{align}
This follows by writing each positive integer prime to $5$ as
$a+5m$ in the series at zero and expanding $(a+5m)^{-2}$.
The right side converges uniformly on the connected affinoid
$|z|\le1$, $|1-z^5|=1$: every $S_j(z^5)$ has integral norm
there, and the $j$th summand has norm at most $5^{-j}$.
The left side is rigid analytic there by
\cite[Theorem~2.7(3)]{BesserDeJeuAlgorithm}; hence the identity
at zero extends to this affinoid.  Here $z\mapsto z^5$ is a
substitution in a function, not a Frobenius automorphism of the
ramified coefficient field $E$.

At $z=-\rho$ one has $z^5=-1$, fixed by all $\sigma_g$.
Thus $e_rF(-\rho)=e_r\operatorname{Li}_2(-\rho)$.
The exact identity
\[
 e_r(\rho^a)=\frac14\omega(a)^r\tau_r\qquad(1\le a\le4)
\]
follows by permuting the four Galois automorphisms.  Applying it
term by term to \eqref{eq:dimtwoDepletion} gives
\[
 \frac{4e_r\operatorname{Li}_2(-\rho)}{\tau_r}
 =\sum_{j\ge0}(-1)^j(j+1)5^j S_j(-1)
       \sum_{a=1}^4\frac{(-1)^a\omega(a)^r}{a^{j+2}}.
\]
Every $S_j(-1)$ is $5$-integral.  Hence the series converges
in $\mathbf Z_5$, and its terms with $j\ge1$ are divisible by $5$.
Its residue is
\[
 \frac12\sum_{a=1}^4(-1)^a a^{r-2}
 =\begin{cases}2,&r=1,\\1,&r=3,\end{cases}
 \quad\text{in }\mathbf F_5.
\]
There is no loss of precision from dividing by $\tau_r$: each
term was divided using the preceding exact identity.

To verify that $\tau_r\ne0$ with the stated valuation, expand
$\rho=1+\pi$.  The coefficient of $\pi^j$ is
$\sum_g\omega(g)^{-r}\binom gj$.  For $j<r$ it is divisible by
$5$, by the power sums of $\mathbf F_5^\times$; for $j=r$ its
residue is $4/r!\ne0$.  Since $v_\pi(5)=4>r$, the unique
lowest-valuation contribution has degree $r$.
\end{proof}

\begin{proposition}[An actual arithmetic coefficient]
\label{prop:dimtwoActualCoefficient}
There is an actual class
$\beta\in K_3(W_{\rm ar})$ of the form
$(1-\sigma_{-1})\beta_0$ such that, for its image in $K_3(W_E)$,
the AMMN point coordinate $b=r_{W_E}(\beta)\in E$ satisfies
\[
 \sigma_{-1}b=-b,\qquad e_1b\ne0,\qquad e_3b\ne0.
\]
The reduction of $\beta$ in $K_3(\mathbf F_5)$ is zero integrally.
\end{proposition}

\begin{proof}
The root-of-unity class $b_2(-\rho)$ is an actual rational
$K_3$-class over the number field $\mathbf Q(\rho)$, in weight two.
It may equivalently be specified by its rational Bloch symbol:
the factor $-\rho$ in its wedge boundary is torsion.
Both $-\rho$ and $1+\rho$ are units in $W_{\rm ar}$, since
$1+\rho\equiv2\pmod\pi$.  Thus this is already a special-unit
symbol, and \cite[Theorem~1.10(2)]{BesserDeJeu} applies directly.
Its local syntomic regulator is a fixed nonzero rational multiple
of $L_{\mathrm{mod},2}(-\rho)=\operatorname{Li}_2(-\rho)$,
because $\log(-\rho)=0$; the sign convention is fixed throughout
\cite[Remark~1.7]{BesserDeJeu}.  Localization at the coefficient
DVR lifts it rationally to $W_{\rm ar}$, since
$K_2(\mathbf F_5)=0$.  Clear one integer denominator to obtain
an actual $\beta_0\in K_3(W_{\rm ar})$, and set
$\beta=(1-\sigma_{-1})\beta_0$.
This is an exact integral antisymmetrization.  Since $\sigma_{-1}$
induces the identity on the residue field, its reduction is zero.

The Chern-compatible syntomic-to-etale comparison at this
ramified coefficient point is a natural Galois-equivariant
isomorphism in twist two.  In the point coordinates it is the
Bloch--Kato exponential; see
\cite[Example~2.2.4 and Propositions~2.2.7, 2.2.9, 2.3.4]{HuberKings}.
Consequently
Lemma~\ref{lem:dimtwoDilogCharacters} proves that both odd
character projections of $c_{2,1}(\beta)$ are nonzero.  The
factor $1-\sigma_{-1}$ multiplies each of them by $2$.

For clarity, no equality between a chosen syntomic coordinate
and the AMMN scalar is needed.  The natural finite-coefficient
Chern classes give a Galois-equivariant $\mathbf Q_5$-linear map
\[
 \pi_3\bigl(K(W_E)^\wedge_5[1/5]\bigr)
       \longrightarrow H^1(E,\mathbf Q_5(2))
\]
through which the Chern class of every actual element factors.
Indeed, pass from the completed spectrum to the inverse limit of
$K_3(E;\mathbf Z/5^\nu)$ and apply its compatible finite Chern
classes.  The groups $H^0(E,\mu_{5^\nu}^{\otimes2})$ are finite,
so their inverse system is Mittag--Leffler and the target inverse
limit is $H^1(E,\mathbf Z_5(2))$.  No vanishing of a derived
inverse limit on the $K$-theory side is required for these
canonical projections.  Naturality, additivity and compatibility
are \cite[Proposition~2.1.1, Lemma~2.3 and Propositions~2.4, 2.8]{WeibelChern}.
The map is $\mathbf Z_5$-linear before rationalization because
each target at level $\nu$ is killed by $5^\nu$.

The point case of the AMMN fiber comparison identifies the source
with $HC_2^{\mathrm{cts}}(W_E/W_E;\mathbf Q_5)=E$, naturally
under $\sigma_g$.  The adjacent positive homotopy groups of the
completed residue-field spectrum vanish by Quillen's finite-field
calculation.  If $e_rb$ were zero, the same projection of its
Chern image would be zero, a contradiction.  This proves the
assertions without assuming that actual $K$-theory surjects onto
any completed vector space.
\end{proof}


% Independent fragment for the exploratory dimension-two draft.
% This file does not modify the main manuscript and makes no K-theory claim.
% Requires the usual lemma/proof environments and amsmath, amssymb.
% Bibliography entry needed: SilvermanAEC, The Arithmetic of Elliptic Curves,
% second edition, GTM 106, Springer, 2009.

\section{The quartic root lattice and a constant elliptic quotient}
\label{sec:quarticellipticgeometry}

Put $k=\overline{\mathbf F}_5$, $W=W(k)$, and $K_0=k(T)$.
Consider the smooth projective quartic over $K_0$
\[
 C:\quad X^4+Y^4+TXY^3=Z^4,
 \qquad C^\circ=C\setminus\{Z=0\}.
\]
For $a\in\mu_4(k)=\mathbf F_5^\times$, write
$P_a=[0:a:1]$ and
$\tau[X:Y:Z]=[-X:-Y:Z]$. Fix $i\in\mu_4(k)$ with $i^2=-1$.
The purpose of this section is entirely geometric: it constructs a
constant elliptic quotient that detects the anti-invariant root
divisors. Survival of its Chern class under Cohen completion is a
separate assertion.

\begin{lemma}[Persistence of the root lattice]
\label{lem:quarticrootlattice}
On the auxiliary open $C^\circ\cap\{Y\ne0\}$ one has
\[
 \operatorname{Pic}(C^\circ\cap\{Y\ne0\})
   =\mathbf Z^4/\mathbf Z(1,1,1,1),
\]
where the four standard generators are the root divisors $P_a$.
In particular, $P_1-P_{-1}$ and $P_i-P_{-i}$ are independent in
$\operatorname{Pic}(C^\circ)$ and in $\operatorname{Jac}(C)(K_0)$.
\end{lemma}
\begin{proof}
Use the regular surface
\[
 D=\operatorname{Spec}
   k[X,U,Y,Y^{-1}]/\bigl(XU-(1-Y^4)\bigr),
 \qquad T=(U-X^3)/Y^3.
\]
It is smooth because the $Y$-derivative $4Y^3$ is invertible.
After inverting $X$, its ring is the Laurent polynomial UFD
$k[X^{\pm1},Y^{\pm1}]$. The four primes over $X$ are
$(X,Y-a)$, with multiplicity one. They generate the divisor class
group. Every relation is the divisor of a unit on $D[1/X]$, hence
of $cX^mY^n$; since $Y$ is already invertible, the only relation is
the sum of the four primes. Thus
$\operatorname{Pic}(D)=\mathbf Z^4/\mathbf Z(1,1,1,1)$.

Every finite fiber $T=t$ is integral, including singular fibers.
Indeed, the change $s=X/Y$, $w=1/Y$ identifies it with
\[
 w^4=s^4+ts+1,\qquad w\ne0.
\]
The polynomial $s^4+ts+1$ is never a square in $k(s)$: a rational
square would be a polynomial square, and comparison with
$(s^2+bs+c)^2$ forces $b=c=0$, contradicting the constant term.
As $k$ contains $\mu_4$, this proves irreducibility of the displayed
fourth-root polynomial. Consequently each removed vertical prime is
the single principal divisor $(T-t)$, of multiplicity one.
Divisor-class localization therefore gives an isomorphism
$\operatorname{Pic}(D)\simeq\operatorname{Pic}(D_{K_0})$.
The generic fiber is the stated auxiliary open of $C$.
Any relation on $C^\circ$, or among the degree-zero classes in its
proper Jacobian, would restrict to a relation in this lattice.
\end{proof}

\begin{lemma}[Two constant elliptic quotients]
\label{lem:quarticconstantquotients}
After a finite separable extension $L/K_0$ there are two finite
degree-two maps
\[
 h_+,h_-:C_L\longrightarrow E_{0,L},
 \qquad E_0:\ V^2=\zeta^3+\zeta.
\]
They satisfy $h_\pm\tau=[-1]h_\pm$. For
$S=\{O,(0,0)\}\subset E_0[2]$, their exact boundary preimages are
\[
 h_\pm^{-1}(S)=\{Z=0\}.
\]
Thus both maps restrict to finite maps $C_L^\circ\to E_0\setminus S$;
no additional points of $C^\circ$ are deleted.
The cover and the maps can be constructed over a finite etale cover
of an open $W$-parameter line, before adjoining any ramified constant.
\end{lemma}
\begin{proof}
On $Y\ne0$ write $s=X/Y$, $w=Z/Y$, so that
$w^4=q_T(s)=s^4+Ts+1$. Its discriminant is
$\Delta=256-27T^4$. On $\operatorname{Spec}W[T,1/\Delta]$ take the
finite etale cover that orders its roots $b_1,\ldots,b_4$. Put
\[
 A=\frac{b_2-b_4}{b_2-b_1},\qquad
 x=A\frac{s-b_1}{s-b_4},\qquad
 \lambda=A\frac{b_3-b_1}{b_3-b_4}.
\]
The quantities $\lambda$, $\lambda-1$, $A$, $A-1$, and
$A-\lambda$ are units. Adjoin a fourth root
$c^4=A(A-1)(A-\lambda)$ and replace $w$ by $cw/(s-b_4)$.
The equation becomes
\[
 w^4=x(x-1)(x-\lambda).
\]
Adjoin $r^2=\lambda$ and $\alpha_\pm^2=r\pm1$. These are again
finite etale extensions: all radicands and the integers $2,4$ are
units. This explicitly constructs the required cover over $W$.

Set
\[
 z=w^2/x,\qquad v_\pm=w(x\pm r)/x.
\]
The identity $z^2=x+r^2/x-1-r^2$ gives
\[
 v_\pm^2=z^3+(r\pm1)^2z.
\]
The substitutions $z=\alpha_\pm^2\zeta$ and
$v_\pm=\alpha_\pm^3V$ give the displayed constant curve $E_0$.
The two maps are the quotients by
\[
 \iota(x,w)=(r^2/x,rw/x),\qquad \tau\iota(x,w)=(r^2/x,-rw/x),
\]
respectively. Their function-field degrees are two: $x$ satisfies
a quadratic equation over the displayed invariant functions, and
$w$ is then recovered rationally. They extend to morphisms of smooth
proper curves. After shrinking the base open, the same formulas give
finite flat maps of degree two over that base.

The four branch points of the fourth-root presentation lie over
$x=0,1,r^2,\infty$. At $0,\infty$ the orders of $(z,v_\pm)$ are
$(-2,-3)$, so both points map to $O$. At $1,r^2$ the orders are
$(2,1)$, so both map to $(0,0)$. Each of these four points is
unramified for either elliptic quotient. Everywhere else $x,w,z$
are finite and nonzero. This proves the exact preimage statement;
the four points in question are precisely the original $Z=0$ boundary.

Finally, the coefficient of $\zeta^4$ in
$(\zeta^3+\zeta)^2$ is $2\ne0$ in $k$, so $E_0$ is ordinary by
the elliptic Hasse-invariant criterion
\cite[Theorem~V.4.1(a)]{SilvermanAEC}. Its automorphism
$[i](\zeta,V)=(-\zeta,iV)$ is the indicated CM action.
\end{proof}

\begin{lemma}[A nonconstant root image and its exact equivariance]
\label{lem:quarticnonconstantroot}
For at least one of $h_+,h_-$, the point $Q=h(P_1)\in E_0(L)$
is non-torsion and nonconstant. For this same map,
\[
 h(P_a)=[a^{-1}]Q\qquad(a\in\mu_4).
\]
In particular, both anti-root differences have non-torsion image,
and their two images have no integral linear relation.
\end{lemma}
\begin{proof}
The three nontrivial involutions form a Klein four group
$G=\{1,\tau,\iota,\tau\iota\}$. Its total quotient is rational,
with coordinate $z=w^2/x$. Let $q:C\to C/\tau$ denote the remaining
elliptic quotient and set
\[
 \Phi=(q_*,h_{+*},h_{-*}),\qquad
 \Psi=q^*+h_+^*+h_-^*.
\]
The divisor pullback--norm formulas give
\[
 \Psi\Phi=(1+\tau)+(1+\iota)+(1+\tau\iota)
          =2+\sum_{g\in G}g=[2].
\]
Here the group norm factors through the zero Jacobian of $C/G$.
Thus $\Phi$ is an isogeny with kernel killed by two. The class
$P_1-P_{-1}$ is non-torsion by Lemma~\ref{lem:quarticrootlattice}
and is killed by $q_*$. At least one of its other two images is
therefore non-torsion. Since $h\tau=-h$, that image is $2Q$.
Every point of $E_0(k)$ is defined over a finite field and is torsion;
thus $Q$ is nonconstant.

The original root rotation $[X:Y:Z]\mapsto[iX:iY:Z]$ sends
$w=Z/Y$ to $-iw$. On either elliptic target this is $[-i]$.
Hence $h(P_i)=[-i]Q$, proving the formula for all four roots.
The two anti-root images are $2Q$ and $-2[i]Q$.
Any nontrivial integral relation would make a nonzero Gaussian
integer annihilate $Q$. Such an endomorphism is an isogeny, which
would make $Q$ torsion, a contradiction.
\end{proof}

The norm relation also keeps track of the diagonal exactly. On the
proper quartic, the rational function $X/Z$ has divisor
$\sum_a P_a-\sum_{Z=0}P$. The boundary computation above therefore
gives, over the smooth parameter open,
\begin{equation}\label{eq:quarticnormdiagonal}
 \operatorname{div}\operatorname{Nm}_h(X/Z)
  =\sum_{a\in\mu_4}\Gamma_{[a^{-1}]Q}
       -2\Gamma_O-2\Gamma_{(0,0)}.
\end{equation}
In particular, the sum of the four norm root classes is zero in
$\operatorname{Pic}((E_0\setminus S)\times U)$.



\section{A scalar Frobenius class at the generic Cohen point}
% Independent insertion for the exploratory dimension-two construction.
% Scalar Frobenius case only; no general generic-point full-faithfulness claim.
\begin{lemma}[A scalar Frobenius class at the generic Cohen point]
\label{lem:scalarFrobeniusCohen}
Let $k=\overline{\mathbf F}_p$, $W=W(k)$, and $K=W[1/p]$.
Let $D_0$ be the coordinate ring of a smooth connected affine curve
$U/k$, and let $D$ be a smooth affine $W$-lift. Write
\[
 \widehat D=\varprojlim_\nu D/p^\nu,
 \qquad C=\varprojlim_\nu S^{-1}D/p^\nu,
 \qquad F=C[1/p],
\]
where $S$ consists of the elements with nonzero reduction in $D_0$.
Assume that an integral Frobenius lift
$\phi:D^\dagger\to D^\dagger$ induces the $q$th-power map on $D_0$,
where $q=p^d>1$, and the Witt Frobenius $\sigma^d$ on $W$.
Let $\omega\in\Omega^1_{D^\dagger/W}[1/p]$ satisfy
\[
 \phi^*\omega-a\omega=dh,
 \qquad a\in W,\quad h\in D^\dagger[1/p].
\]
If the class of $\omega$ in $H^1_{\mathrm{rig}}(U/K)$ is nonzero,
then its image in $\Omega^1_{C/W}[1/p]/dF$ is nonzero.
Here the differentials of $C$ are continuous differentials.
\end{lemma}
\begin{proof}
We first recall that the constant field of the continuous differential
field $F/K$ is exactly $K$. On an etale-coordinate open of $U$, the
integral Taylor homomorphism gives Hasse derivatives $D_j$ on $C$,
with $j!D_j=\partial^j$. Thus $df=0$ implies that all positive Hasse
derivatives of the reduction of an integral $f$ vanish. Their common
kernel in $k(U)$ is its maximal perfect subfield, namely $k$.
Subtract a Witt-vector lift of this constant and divide by $p$;
iteration and completeness show that $f\in W$. Scaling treats $F$.

Suppose now that $\omega=df$ in $F$. The displayed Frobenius relation
gives
\[
 \phi(f)-af=h+c\qquad(c\in K).
\]
Multiplying $f,h,c,\omega$ by one common power of $p$ makes all of
$f\in C$, $h\in\widehat D$, and $c\in W$ integral. The lift $\phi$
extends to $C$, fixes $p$, and preserves $\widehat D$.

More generally, if $f_n\in C$ and
$\phi(f_n)-af_n=r_n\in\widehat D$, its reduction satisfies
\[
 \overline f_n^{\,q}-\overline a\,\overline f_n
       =\overline r_n\in D_0.
\]
At a closed point of $U$ where $\overline f_n$ had a pole of order
$m>0$, the first term would have pole order $qm$, whereas the second
has order at most $m$. Cancellation is impossible. Hence
$\overline f_n\in D_0$. Choose a lift $g_n\in D$, and put
\[
 f_{n+1}=(f_n-g_n)/p,\qquad
 r_{n+1}=(r_n-\phi(g_n)+ag_n)/p.
\]
The natural inclusion $D_0\hookrightarrow k(U)$ implies
$\widehat D\cap pC=p\widehat D$, so $r_{n+1}\in\widehat D$.
The same relation holds for $n+1$. Starting with $r_0=h+c$, this
constructs a convergent expansion $f=\sum_{n\geq0}p^ng_n$ in
$\widehat D$. Therefore $\omega$ is exact already on the convergent
lift of $U$.

The injectivity
$H^1_{\mathrm{rig}}(U/K)\to H^1_{\mathrm{conv}}(U/K)$
of \cite[Theorem~3.2.1]{CaroDAddezioInjectivity} contradicts the
assumed nonzero rigid class. This proves the lemma.
\end{proof}

% The application uses a = alpha^d in Z_p^times and the pullback of the
% unit-root line of a constant ordinary elliptic curve. The proof only
% needs a integral; do not generalize it to arbitrary matrix coefficients
% or unspecified isocrystals without further argument.


% Proposed standalone dimension-two draft: internal review only.
% Do not insert into the existing main manuscript without a new review.
% Input dimension_two_scalar_frobenius_cohen_insertion.tex before this file.
% Geometry is supplied by dimension_two_quartic_prym_root_geometry_audit.md.

\section{An elliptic detector after generic Cohen completion}
\label{sec:dimtwoEllipticDetector}

Put $k=\overline{\mathbf F}_5$, $W=W(k)$, and $K=W[1/5]$.
Let $E_0/k$ be the ordinary elliptic curve
\[
 E_0:\quad v^2=z^3+z,
 \qquad S=\{O,(0,0)\},\qquad E_0^\circ=E_0\setminus S.
\]
The equation and the two deleted sections have their evident smooth
lifts over $W$; write $A$ for the coordinate ring of this affine lift
and $\widehat A$ for its $5$-adic completion.

\begin{lemma}[A bounded functional in the elliptic direction]
\label{lem:ellipticContinuousFunctional}
There is a bounded $K$-linear functional
\[
 \ell:\Omega^1_{\widehat A/W}[1/5]\longrightarrow K
\]
that kills $d\widehat A[1/5]$, is nonzero on the image of the proper
slope-one line, and is zero on the image of the proper slope-zero
line.  A nonzero scalar multiple of
$\ell$ sends $\Omega^1_{\widehat A/W}$ into $W$.
\end{lemma}

\begin{proof}
Ordinarity follows from the nonzero coefficient $2$ of $z^4$ in
$(z^3+z)^2$ modulo $5$.  Proper degree-one cohomology injects into
$H^1_{\mathrm{rig}}(E_0^\circ/K)$ by the Gysin sequence.  The slope-one
part maps isomorphically to the corresponding convergent part by
\cite[Theorem~7.3.1]{CaroDAddezioInjectivity}, and its image has zero
intersection with the closure of zero by
\cite[Theorem~7.3.6]{CaroDAddezioInjectivity}.  The slope-zero part is
in that closure by \cite[Theorem~7.2.1]{CaroDAddezioInjectivity}.

Let $M=\Omega^1_{\widehat A/W}[1/5]$ and
$N=\overline{d\widehat A[1/5]}\subset M$ for the affinoid topology.
The nonzero proper slope-one line therefore has nonzero image in the
Hausdorff Banach space $M/N$.  A continuous linear functional on this
quotient can be chosen nonzero on that line.  Indeed, a complete
discretely valued field is spherically complete, so the nonarchimedean
Hahn--Banach theorem applies; equivalently one can use an orthonormal
basis of this discretely normed Banach quotient.  Pulling the functional
back to $M$ gives $\ell$.  Boundedness gives the last assertion.
\end{proof}

Let $U/k$ be a smooth connected affine curve, with a smooth affine
$W$-lift $D$.  Write $C$ for the $5$-adic completion of $D$ after
inverting all elements with nonzero reduction, and write $F=C[1/5]$.
Thus $C$ is the generic Cohen ring occurring in
Lemma~\ref{lem:scalarFrobeniusCohen}.  All differentials of $C$ are
continuous differentials over $W$.

\begin{lemma}[The completed tensor operator]
\label{lem:ellipticCompletedTensor}
For
\[
 B=(\widehat A\otimes_W C)^{\wedge}_5,
\]
the functional in Lemma~\ref{lem:ellipticContinuousFunctional} defines
a degree-minus-one map of complexes
\[
 L:\Omega^\bullet_{B/W}[1/5]
       \longrightarrow \Omega^\bullet_{C/W}[1/5][-1].
\]
On a product of an elliptic one-form $\alpha$ and a base one-form
$\gamma$, it is $L(\alpha\wedge\gamma)=\ell(\alpha)\gamma$.
\end{lemma}

\begin{proof}
Scale $\ell$ so that it preserves the integral lattices.  For each
$\nu$, tensor
\[
 \Omega^1_{\widehat A/W}/5^\nu\longrightarrow W/5^\nu
\]
with $C/5^\nu$ and take the inverse limit.  This defines the required
bounded operator on the completed tensor product.  Tensoring instead
with $\Omega^1_{C/W}/5^\nu$ defines its degree-two part.  The usual
formula for differentials of a tensor product, applied at each finite
precision, gives
\[
 \Omega^1_{B/W}
  =\bigl(\Omega^1_{\widehat A/W}\widehat\otimes_W C\bigr)
    \oplus\bigl(\widehat A\widehat\otimes_W\Omega^1_{C/W}\bigr),
 \qquad
 \Omega^2_{B/W}
  =\Omega^1_{\widehat A/W}\widehat\otimes_W\Omega^1_{C/W}.
\]
Both factors have relative dimension one.  In an etale coordinate
$T$ on the base, write a one-form as $\alpha+b\,dT$, where $\alpha$
is an elliptic one-form with base-dependent coefficients.  Then
\[
 L_2\bigl(d(\alpha+b\,dT)\bigr)
   =-d_C\bigl(\ell(\alpha)\bigr),
\]
because $\ell(d_A b)=0$ and the fixed bounded $W$-linear operator
$\ell$ commutes with the continuous derivation in $T$.  This is
precisely the chain-map identity for the displayed shift.  The
construction is intrinsic and hence agrees on different etale
coordinates.  In particular, exact two-forms have exact images.

This construction takes inverse limits of explicitly compatible
operators on forms.  It does not interchange cohomology with an
infinite localization or with a completed tensor product.
\end{proof}

\begin{proposition}[A nonconstant elliptic point]
\label{prop:ellipticCohenDetector}
Suppose $Q:U\to E_0^\circ$ is nonconstant and the finite-type data
are lifted over $W$, after shrinking $U$ if necessary.  Let
$\mathcal L_Q$ be the line bundle on the completed product $B$
obtained from the divisor $\Gamma_Q-\Gamma_{-Q}$.  Then
\[
 c_1(\mathcal L_Q)\ne0
 \quad\text{in}\quad
 \Omega^2_{B/W}[1/5]\big/d\Omega^1_{B/W}[1/5].
\]
The same conclusion holds after any finite constant coefficient
extension of $W$.
\end{proposition}

\begin{proof}
Normalize the Poincare bundle on $E_0\times E_0$ at the origin.
The line bundle of $\Gamma_Q-\Gamma_O$ agrees with its pullback
along $\operatorname{id}\times Q$ up to a line bundle from the base.
In the difference between $Q$ and $-Q$ these base factors cancel,
since $[-1]^*\mathcal O_{E_0}(O)=\mathcal O_{E_0}(O)$.

Choose dual generators $\eta$ and $\omega$ on the proper slope-zero
and slope-one lines of $H^1(E_0)$, respectively.  The Poincare pairing
and the rigid Kunneth formula identify the mixed part of the Chern
class with
\[
 \pm2\bigl(\omega\otimes Q^*\eta
                   -\eta\otimes Q^*\omega\bigr).
\]
Only this mixed part occurs: the two graph divisors have the same
relative degree, and the base is an affine curve.  The contraction
of Lemma~\ref{lem:ellipticCompletedTensor} consequently sends this
class to
\[
 \pm2\ell(\omega)\,Q^*\eta.
\]
This computation is made in finite-dimensional rigid cohomology of
the finite-type source.  Choose a finite sum of products of dagger
one-forms representing that Kunneth class.  Its difference from a
Chern representative is exact in the dagger complex and remains
exact after passage to the completed complex and application of $L$.
Thus no finite-dimensional description of the completed target is
being assumed.

The rational map $Q$ extends to the smooth proper compactification
of $U$ and is finite.  Its pullback on proper rational crystalline
$H^1$ is injective, since trace composed with pullback is its nonzero
degree; this also covers an inseparable map.  Restriction to $U$ is
injective in rigid degree one.  Therefore $Q^*\eta\ne0$ in
$H^1_{\mathrm{rig}}(U/K)$.

All special-fiber data descend to a finite field $\mathbf F_{5^d}$.
The smooth affine lift of the base curve admits an integral
Frobenius lift on its weak completion, by
\cite[Introduction, p.~198]{DavisLangerZink}; iterate it $d$ times.
The induced action on rigid cohomology is the canonical one, so the
scalar class below satisfies the exact-form hypothesis of
Lemma~\ref{lem:scalarFrobeniusCohen}.
The unit-root Weil number of $E_0/\mathbf F_5$ is
$\alpha=1-2i\in\mathbf Z_5^\times$, with $i\equiv2\pmod5$:
the characteristic polynomial is $X^2-2X+5$.  Hence $Q^*\eta$ is a
scalar Frobenius class with multiplier $\alpha^d$.
Lemma~\ref{lem:scalarFrobeniusCohen} proves that its image in
$H^1_{\mathrm{dR}}(F/K)$ is still nonzero.  Since
$\ell(\omega)\ne0$, the contraction is nonzero.  Exactness of the
original two-form would contradict the chain-map identity.

A finite coefficient extension is flat and commutes with the kernels
and images of the complexes in question, so it preserves this
nonzero class.
\end{proof}

\begin{lemma}[Tame trace and the norm]
\label{lem:tameQuarticTrace}
Let $h:\mathcal C^\circ\to E_0^\circ\times U$ be either of the
degree-two quartic quotient maps, with involution $\iota$.  On the
completed generic Cohen algebras there is a trace on total continuous
de Rham complexes, relative to $W$, satisfying
\[
 \operatorname{Tr}\circ h^*=2,\qquad
 \operatorname{Tr}(c_1(\mathcal L))
      =c_1(\operatorname{Nm}_h\mathcal L).
\]
In particular, a nonzero Chern class of the norm implies a nonzero
Chern class on the source.
\end{lemma}

\begin{proof}
The invariant differentials are the differentials on the quotient.
This is etale descent away from ramification; at a ramification point
the relative parameters have the form $s=t^2$, and an invariant
one-form is a multiple of $t\,dt=ds/2$.  The same calculation includes
the base differential $dT$ and its wedge products.  Thus summing
$1+\iota^*$ and descending gives a strict chain map, integral because
$2$ is a unit.  It is bounded and therefore extends through the
completion.  The first identity is immediate.

For a line bundle on a degree-two quotient,
$h^*\operatorname{Nm}_h\mathcal L
 \simeq\mathcal L\otimes\iota^*\mathcal L$; the divisor identity
also holds at ramification points, where it counts multiplicities.
Pulling back both sides of the claimed Chern identity gives the same
class.  The first trace identity makes $h^*$ injective after inverting
$5$, proving the second identity.  Equivalently it follows directly
from $d\log\operatorname{Nm}=\operatorname{Tr}\,d\log$.
\end{proof}

\subsection{Replacing the scaled lift while preserving its root lines}

Let
\[
 V=\widehat{W[T]_{(5)}}_5,
 \quad W'=W[\rho_5],\quad u=\rho_5-1,
 \quad V'=V\otimes_W W'.
\]
The finite algebra $V'$ is already complete and is a DVR with
uniformizer $u$.  Restricted power series are denoted by angle brackets.

\begin{lemma}[Formal replacement]
\label{lem:quarticFormalReplacement}
There is a $V'$-linear isomorphism, reducing to the identity modulo $u$,
between
\[
 R_{\mathrm{good}}=
 V'\langle X,Y\rangle/
 \left(X^4+\frac{\Phi_5(1+uY)}{u^4}+TXY^3\right)
\]
and
\[
 R_{\mathrm{cyc}}=
 V'\langle X,Y\rangle/(X^4+Y^4+TXY^3-1).
\]
It carries the Picard class of each root ideal
$\bigl(X,Y-(\rho_5^g-1)/u\bigr)$ to the class of
$\bigl(X,Y-[g]\bigr)$, where $g\in\mathbf F_5^\times$ and $[g]$
is its Teichmuller fourth root of unity.
\end{lemma}

\begin{proof}
The equation $\Phi_5(1+u)=0$ gives $5/u^4\equiv-1\pmod u$.
Every other lower coefficient of $\Phi_5(1+uY)/u^4$ is divisible by
$u$.  Both reductions are therefore $G(X,Y)=1$, with
$G=X^4+Y^4+TXY^3$.  This affine curve is smooth: simultaneous
vanishing of both partial derivatives would give $4G=0$ by Euler's
identity, contradicting $G=1$.  The two lifts are $u$-torsion free
(their equations are monic in $Y$), complete, and formally smooth
over $V'$.

Formal smoothness successively lifts the identity of the reductions
to compatible maps modulo $u^\nu$.  Completeness gives a $V'$-linear
map between the displayed rings.  It is surjective by successive
approximation.  A kernel element is successively divisible by every
$u^\nu$, using torsion-freeness of the target; separatedness therefore
makes the kernel zero.

The root ideals are invertible: along each root the derivative in
$Y$ is a unit, and their quotients are $V'$.  Their reductions are
the same ideals $(X,Y-g)$.  Finally $u$ belongs to the Jacobson
radical, and $\operatorname{Pic}(R_{\mathrm{cyc}})\to
\operatorname{Pic}(R_{\mathrm{cyc}}/u)$ is injective.  Indeed, an
isomorphism of two reduced projective lines lifts to a homomorphism,
and Nakayama at every maximal ideal makes it an isomorphism.  This
identifies the transported root classes.

Because the isomorphism is $V'$-linear, it preserves all arithmetic
coefficient classes and the parameter $T$.  The induced isomorphism
of continuous complexes is taken relative to $W'$, and in particular
retains $dT$.
\end{proof}

\begin{corollary}[A visible anti-root Chern class]
\label{cor:quarticRootChernDetected}
For the cyclic model in Lemma~\ref{lem:quarticFormalReplacement},
the map sending the anti-root Picard lattice to ordinary continuous
degree-two de Rham cohomology over $W'$ is nonzero.  The same holds
for the corresponding root lines on $R_{\mathrm{good}}$.
\end{corollary}

\begin{proof}
Use the explicit quartic geometry after a finite etale cover of the
parameter curve, constructed first over $W$ and then tensored with
$W'$.  It supplies a degree-two map $h$ to $E_0^\circ$ for which
$Q=h(P_1)$ is nonconstant and $h(P_{-1})=-Q$.  The norm of
$\mathcal O(P_1-P_{-1})$ is consequently
$\mathcal O(\Gamma_Q-\Gamma_{-Q})$.
Proposition~\ref{prop:ellipticCohenDetector} and
Lemma~\ref{lem:tameQuarticTrace} prove that its source Chern class is
nonzero after generic Cohen completion and finite coefficient
extension.  Nonzero pullback implies nonzero before the finite
parameter cover.  Formal replacement proves the last assertion.
This statement does not assert linear independence after extending
the entire root lattice to $\mathbf Q_5$; only a nonzero map is needed.
\end{proof}

% Bibliography entry for the standalone draft:
% \bibitem{CaroDAddezioInjectivity}
% D.~Caro and M.~D'Addezio,
% \emph{Injectivity failure in crystalline comparisons},
% arXiv:2511.11444 (2025).
% \url{https://arxiv.org/abs/2511.11444}.
% The numbered inputs used here are Theorems 3.2.1, 7.2.1, 7.3.1, 7.3.6.


% Independent assembly for the exploratory dimension-two paper.
% No edits to the main manuscript.  The geometric detector and the explicit
% coefficient lemma must be supplied separately by the parent manuscript.

\section{From the completed detector to an actual generic-kernel class}
\label{sec:dimensiontwoactualassembly}

Put $k=\overline{\mathbf F}_5$, $W_0=W(k)$, and
\[
 V_{\mathrm{ar}}=\mathbf Z_{(5)}[T]_{(5)},\qquad
 V=\widehat{V_{\mathrm{ar}}}^{\,(5)}
   =\widehat{\mathbf Z_5[T]_{(5)}}^{\,(5)},\qquad
 V_{\mathrm{big}}=\widehat{W_0[T]_{(5)}}^{\,(5)}.
\]
Our two source rings are
\begin{align}
 A_2&=\left(\mathbf Z_{(5)}[T,x,y]/
       \bigl(\Phi_5(1+y)+x^4+Txy^3\bigr)\right)_{(5,x,y)},
       \label{eq:dimensiontwoArithmeticSource}\\
 \widehat A_2&=\widehat{A_2}^{\,(x,y)}
      \simeq V[[x,y]]/
       \bigl(\Phi_5(1+y)+x^4+Txy^3\bigr).
       \label{eq:dimensiontwoCompletedSource}
\end{align}
The first is essentially of finite type over $\mathbf Z$; the
second is its maximal-ideal completion.
Let $\rho$ be a primitive fifth root of unity,
$\pi=\rho-1$, $E=\mathbf Q_5(\rho)$,
$W_{\mathrm{ar}}=\mathbf Z_{(5)}[\rho]$,
$W_E=\mathbf Z_5[\rho]$, $W'=W_0[\rho]$, and $V'=V[\rho]$.
Write $\sigma_g(\rho)=\rho^g$ for $g\in\mathbf F_5^\times$.
For every ring $A$, the notation
\[
 \mathcal K(A)=K(A)^{\wedge}_5[1/5]
\]
means completion of the \emph{spectrum} followed by rationalization;
it does not mean $K(A)\otimes\mathbf Q_5$ or completion of the ring.
The same convention applies to relative $K$-theory. The primitive
denominators in $\mathbf Z_5[T]$ remain nonzero after reduction in
$k[T]$, so there is a continuous local injection
$V\hookrightarrow V_{\mathrm{big}}$. Both are discrete valuation
rings with uniformizer $5$, so this injection is flat. Both source
rings have residue field $\mathbf F_5(T)$; the algebraic closure
is used only in the test algebra. The displayed description of
$\widehat A_2$ follows by first completing the regular local
ambient ring $V_{\mathrm{ar}}[x,y]_{(5,x,y)}$, whose completion
is $V[[x,y]]$, and then taking its hypersurface quotient.

We use a finite unramified extension $\widetilde V/V_{\mathrm{big}}$ obtained from a
finite etale cover of an open of the $T$-line, chosen over $W_0$.
Set $\widetilde V'=\widetilde V[\rho]$. The test algebra is
\begin{equation}
 B=\widetilde V'\langle X,Y\rangle/
   \left(X^4+\frac{\Phi_5(1+\pi Y)}{\pi^4}+TXY^3\right),
 \qquad C=B/\pi.
\label{eq:dimensiontwotestalgebra}
\end{equation}
The polynomial displayed in this formula has integral coefficients.
Its reduction is $X^4+Y^4+TXY^3-1$. Euler's identity, and the fact
that $4$ is invertible, show that this affine curve is smooth.
Consequently $B$ is a regular, $\pi$-torsion-free, complete ring,
and $C$ is regular. Both rings have finite Krull dimension.

\begin{lemma}[The precise continuous cyclic input]
\label{lem:dimensiontwoAMMNinput}
The ring $B$ has continuous relative differentials over $W'$ which
are finite projective of rank two, and admits a natural relative
character
\[
 \chi_B:\pi_3\mathcal K(B,(\pi))
       \longrightarrow HC^{\mathrm{cts}}_2(B/W';\mathbf Q_5).
\]
Under the normalized mixed HKR comparison the latter has
$\Omega^2_{B/W'}[1/5]/d\Omega^1_{B/W'}[1/5]$ as its top-form
summand. The character is compatible with multiplication by
$K_0(B)$ through the negative cyclic Chern character.
\end{lemma}
\begin{proof}
The required scope is that $B$ is the $5$-adic completion of an
essentially smooth $W'$-algebra of relative dimension two. The
unramified parameter cover is first constructed on a smooth
$W_0$-curve and then localized at its generic special point.
Adjoining $\rho$ and the displayed affine equation gives such a
model near the entire special fiber. If necessary, localize the
model at one element $j\equiv1\pmod\pi$ contained in its smooth
locus; this leaves its completion unchanged. This also verifies
the stated rank of continuous differentials. In particular,
the base for the differential complex is $W'$, so the parameter
differential is retained.

Apply the henselian fiber square of \cite[Theorem~A]{AMMN}, its
affine comparison \cite[Proposition~4.2]{AMMN}, and the construction
in the proof of \cite[Proposition~4.10]{AMMN}. Here is why their
relative Hochschild step applies to this essentially smooth model.
Writing $W'=W_0[\pi]$ and using the Eisenstein polynomial of
$\pi$, the completed complex $L_{W'/\mathbf Z}$ is represented
by $[W'\xrightarrow{f'(\pi)}W']$; one fixed power of $5$
annihilates this complex. The transitivity triangle and the finite
projective relative cotangent module imply that, in each fixed
degree range, the cone of the absolute-to-$W'$-relative Hochschild
comparison has bounded $5$-power torsion. Only finitely many HKR
graded pieces contribute in that range: the $j$th piece is
$j$-connective. Derived completion and circle orbits preserve this
fixed-range argument. For circle orbits one uses their skeletal
filtration, whose omitted cells become arbitrarily connective;
one does not interchange an unrestricted inverse limit with
homotopy orbits.

Finally the normalized mixed HKR map in relative degrees $0,1,2$
uses only the denominator $2$, a unit at $5$. It identifies the
completed relative mixed complex with the finite complex of
continuous differential forms. The comparison of the negative
and periodic cyclic fiber squares therefore gives the asserted
ordinary cyclic target. To check its module structure, put
$H=\operatorname{THH}(B)$ in cyclotomic spectra. The intermediate
arrows in the fiber-square construction are $H$-linear arrows
between objects $H\otimes D$. Lax monoidality of $\operatorname{TC}$
sends these to modules over $\operatorname{TC}(B)$, and the
equivalences used in the construction are inverted in that module
category. The projection to relative Hochschild homology is
$H$-linear as well. The circle norm fiber has the usual negative
cyclic module action, by the norm projection formula, including
its connecting map. Restriction along the multiplicative
cyclotomic trace gives the claimed $K(B)$-action; in mixed complexes
it is the shuffle-plus-cyclic-shuffle action. Thus this assertion
uses the module construction, rather than naturality alone.
Appendix~\ref{sec:dimtwoAMMNAppendix} supplies the full fixed-range argument for this essentially smooth extension.
\end{proof}

For $g\in\mathbf F_5^\times$ put
\[
 a_g=\frac{\rho^g-1}{\pi}\in\widetilde V',\qquad
 J_g=(X,Y-a_g)\subset B,\qquad
 \eta_g=1-[J_g]\in K_0(B).
\]
These ideals are invertible: on their support the $Y$-derivative
of the equation is a unit, and $X$ is a local equation.
Set $\gamma_g=-c_1^{\mathrm{dR}}(J_g)$, with the normalized
negative cyclic Chern-character convention.

Let
\[
 r_{W_E}:K_3(W_E)\longrightarrow E
\]
be the point character obtained from the same fiber square.
Its definition uses
$\pi_3\mathcal K(\mathbf F_5)=
 \pi_4\mathcal K(\mathbf F_5)=0$; the relative point class has
ordinary cyclic representative $t^{-1}b$ for $b=r_{W_E}(\beta)$,
where $|t|=-2$. This character is natural under coefficient maps
and under the automorphisms $\sigma_g$.

\begin{proposition}[Actual assembly]
\label{prop:dimensiontwoactualassembly}
Let $\beta\in K_3(W_{\mathrm{ar}})$ be an actual class with
$\sigma_{-1}\beta=-\beta$ and with zero reduction in
$K_3(\mathbf F_5)$. Suppose the geometric detector proves
\begin{equation}
 \sum_{g\in\mathbf F_5^\times}\sigma_g(b)\gamma_g\ne0
 \quad\text{in}\quad
 \Omega^2_{B/W'}[1/5]/d\Omega^1_{B/W'}[1/5],
 \qquad b=r_{W_E}(\beta|_{W_E}).
\label{eq:dimensiontwodetectedcoefficient}
\end{equation}
Then there is an actual class $c\in K_3(A_2)$ whose image
$\widehat c\in K_3(\widehat A_2)$ has infinite order, whereas
\[
 c|_{\operatorname{Frac}A_2}=0,\qquad
 \widehat c|_{\operatorname{Frac}\widehat A_2}=0
\]
integrally. In particular the rational generic map in degree
three is not injective for either regular local ring of
dimension two and residue field $\mathbf F_5(T)$.
The reduction hypothesis holds automatically for a coefficient
of the form $(1-\sigma_{-1})\beta_0$.
\end{proposition}
\begin{proof}
The last assertion follows because $\sigma_{-1}$ acts trivially
on the residue field. The even group $K_4(\mathbf F_5)$ is zero,
so the image of the specified coefficient has a unique actual
lift to $K_3(W_E,(\pi))$. Pull it back to the coefficient extension and
multiply by the line-bundle classes to form
\[
 \delta_B=\sum_g\sigma_g(\beta)\eta_g\in K_3(B),
\]
together with a chosen actual relative lift obtained in this way.
There is no surjectivity assumption from an actual $K$-group to
its spectral completion.

By Lemma~\ref{lem:dimensiontwoAMMNinput}, the image of this
relative lift is the sum of the products of
$t^{-1}\sigma_g(b)$ with $\operatorname{ch}^-(\eta_g)$.
The latter has rank zero and degree-two term $t\gamma_g$.
All its higher differential-form terms vanish in relative
dimension two. In the shuffle-plus-cyclic-shuffle product, the
cyclic-shuffle corrections put the scalar $\sigma_g(b)$ into a
normalized bar slot relative to $W'$, and hence vanish.
Thus the top-form component is exactly
$\sum_g\sigma_g(b)\gamma_g$, up to a common nonzero factor if a
different Chern-character normalization is used. Hypothesis
\eqref{eq:dimensiontwodetectedcoefficient} makes the relative
completed class nonzero.

We now justify both passages out of this relative target.
The special fiber $C$ is the filtered colimit, over affine opens
of the smooth parameter curve, of smooth affine surfaces over
the perfect field $k$. Geisser--Levine
\cite[Theorem~8.4]{GeisserLevineCharP} gives
\[
 K_j(C;\mathbf Z/5^\nu)=0\qquad(j>2,\ \nu\geq1).
\]
Finite-coefficient $K$-theory commutes with these filtered
colimits. The Milnor sequence for spectral completion uses
these groups in degrees $j$ and $j+1$, so
\[
 \pi_j\mathcal K(C)=0\qquad(j\geq3).
\]
In the relative fiber sequence the vanishing of
$\pi_4\mathcal K(C)$ therefore gives an injection
\begin{equation}
 \pi_3\mathcal K(B,(\pi))\hookrightarrow\pi_3\mathcal K(B).
\label{eq:dimensiontworelativeinjection}
\end{equation}
Since $B$ and $C$ are regular and $\pi$ is a non-zero-divisor,
localization and devissage identify the spectrum of perfect
complexes supported on $V(\pi)$ with $K(C)$. Completion and
rationalization are exact functors of spectra. The vanishing
of $\pi_3\mathcal K(C)$ gives the second injection
\begin{equation}
 \pi_3\mathcal K(B)\hookrightarrow\pi_3\mathcal K(B[1/\pi]).
\label{eq:dimensiontwogenericinjection}
\end{equation}
Thus the image of the actual class $\delta_B$ in the last group
is nonzero. No assertion that the actual group $K_3(C)$ is
torsion has been made or used.

It remains to match this class to the arithmetic source and its
completion. In the regular local ambient ring
$V_{\mathrm{ar}}[x,y]_{(5,x,y)}$ of dimension three, the defining
equation is congruent to $5$ modulo $(5,x,y)^2$. Thus $A_2$ is
regular local of dimension two and has maximal ideal $(x,y)$.
The same is true of its completion. Its Cartier divisor
\[
 i:\operatorname{Spec}D_{\mathrm{ar}}\hookrightarrow
       \operatorname{Spec}A_2,
 \qquad
 D_{\mathrm{ar}}=A_2/(x)
       \simeq W_{\mathrm{ar}}[T]_{(\pi)}
\]
is identified by $y\mapsto\pi$. Pull $\beta$ to
$K_3(D_{\mathrm{ar}})$ and define the actual class
$c=i_*(\beta)\in K_3(A_2)$, using regularity to identify
$G$- and $K$-theory. Localization at $x$ shows that its image
in $K_3(\operatorname{Frac}A_2)$ is zero integrally.
The completion map is flat, so its image $\widehat c$ is the
Gysin class of the same coefficient on
$\widehat A_2/(x)\simeq V'$. In particular its generic image
also vanishes integrally.

First make the \emph{finite flat coefficient extension}
$\mathbf Z_{(5)}\to W_{\mathrm{ar}}$,
and then invert $5$. On
\[
 S=(A_2\otimes_{\mathbf Z_{(5)}}W_{\mathrm{ar}})[1/5]
\]
the divisor $x=0$ splits into four disjoint components
$y=\rho^g-1$. Write $I_g=(x,y-(\rho^g-1))\subset S$ for
their invertible ideals. Flat base change for $i_*$, followed
by the length-one resolutions of these Cartier divisors, gives
\[
 c|_S=\sum_g\sigma_g(\beta)(1-[I_g])\quad\text{in }K_3(S).
\]
Here the coefficient on the component $y=\rho^g-1$ is
$\sigma_g(\beta)$ because its copy of $1+y$ maps to $\rho^g$.

We spell out the denominator and completion maps. Under the
substitution $x=\pi X$, $y=\pi Y$, every denominator
$s\in\mathbf Z_{(5)}[T,x,y]\setminus(5,x,y)$ reduces modulo
$\pi$ to the nonzero polynomial $\overline s(T,0,0)$ in the
function field of the parameter cover. Its image is therefore
a unit in $C$, hence a unit in the $\pi$-adically complete ring
$B$. The equation of $A_2$ maps to $\pi^4$ times the equation
of $B$. This defines $A_2\to B$, and the coefficient inclusion
$W_{\mathrm{ar}}\to\widetilde V'$ gives
$S\to B[1/\pi]$. Moreover $(x,y)B\subset\pi B$, so the map
extends uniquely and continuously to
\[
 A_2\longrightarrow\widehat A_2\longrightarrow B.
\]
Concretely, coefficients in $V_{\mathrm{ar}}$ extend to $V$,
and substituting $\pi X,\pi Y$ into a formal series converges
$\pi$-adically. No assertion that a $K$-group commutes with this
completion is used.

Pull back the preceding equality as an equality of actual
$K$-classes. The line $I_g$ pulls back to $J_g[1/\pi]$.
This can be checked on its local Cartier equation; its image is
generated by $\pi X$ and $\pi(Y-a_g)$, and $X$ remains a
non-zero-divisor. Consequently
\[
 c|_{B[1/\pi]}=\widehat c|_{B[1/\pi]}
       =\delta_B|_{B[1/\pi]}.
\]
This step pulls back line bundles after the flat coefficient
calculation. It does not assert that $A_2\to B$ or
$\widehat A_2\to B$ is flat, or that the two constructions of
the class agree in $K_3(B)$ before inverting $\pi$.
By \eqref{eq:dimensiontwogenericinjection} the last actual class
has nonzero image in a rationalized completed spectrum. It
therefore has infinite order. Both $c$ and $\widehat c$ have
infinite order, as required.
\end{proof}

\begin{remark}[Scope of the geometric and coefficient inputs]
The preceding proposition is conditional only on the explicit
ordinary de Rham nonvanishing in
\eqref{eq:dimensiontwodetectedcoefficient}. The cyclic formal
replacement transports the four lines to the standard root
lines because their reductions agree and $\operatorname{Pic}(B)
\to\operatorname{Pic}(C)$ is injective. The constant ordinary
elliptic quotient, its bounded degree-minus-one map, and
Lemma~\ref{lem:scalarFrobeniusCohen} prove the needed nonzero
root-difference class after all Cohen denominators. Proposition~\ref{prop:dimtwoActualCoefficient} supplies the actual coefficient, and Proposition~\ref{prop:dimtwoWeightedNonzero} checks the displayed weighted nonvanishing. An element manufactured by a $\mathbf Q_5$-linear
idempotent on completed $K$-theory is not a substitute for
that actual coefficient.
\end{remark}

% Bibliography keys needed:
% AMMN: Antieau--Mathew--Morrow--Nikolaus, On the Beilinson fiber square.
% GeisserLevineCharP: Geisser--Levine, The K-theory of fields in
% characteristic p, Invent. Math. 139 (2000), 459--493, Theorem 8.4.
% CaroDAddezioInjectivity: Caro--D'Addezio, Injectivity failure in
% crystalline comparisons, arXiv:2511.11444, Theorem 3.2.1.


% Root-written link between the arithmetic and geometric inputs.
\section{The weighted root class and the conclusion}
\label{sec:dimtwoWeightedConclusion}

Continue with the finite parameter cover of
Lemma~\ref{lem:quarticconstantquotients}, and choose the quotient
$h$ of Lemma~\ref{lem:quarticnonconstantroot}.  Write $Q=h(P_1)$.
Let $\eta$ be a generator of the proper unit-root line of $E_0$,
and let $\chi_0:\mathbf F_5^\times\to\mathbf Z_5^\times$ be the
character through which $[g]^*$ acts on this line.  Put
$e_{\chi_0}=\frac14\sum_g\chi_0(g)^{-1}\sigma_g$.  The character $\chi_0$ is one of
the two odd characters.  More precisely, for the lift of $i$ with
$i\equiv2\pmod5$, $[i]^*\eta=-i\eta$, so $\chi_0=\omega^3$.
Indeed, the holomorphic differential of $v^2=z^3+z$ has character
$i$, and Frobenius is the CM endomorphism $1+2i$; its action on the
other line is the unit $1-2i$.

\begin{proposition}[The actual coefficient is seen by the detector]
\label{prop:dimtwoWeightedNonzero}
Let $b$ be the AMMN point coordinate of the actual coefficient
from Proposition~\ref{prop:dimtwoActualCoefficient}.  On the scaled
test algebra of \eqref{eq:dimensiontwotestalgebra}, put
$\gamma_g=-c_1(J_g)$.  Then
\[
 \Xi=\sum_{g\in\mathbf F_5^\times}\sigma_g(b)\gamma_g
 \ne0
 \quad\text{in}\quad
 \Omega^2_{B/W'}[1/5]/d\Omega^1_{B/W'}[1/5].
\]
\end{proposition}

\begin{proof}
Formal replacement transports each root line to the cyclic
quartic root line, while fixing $W'$ and the parameter.  Apply
the tame trace of Lemma~\ref{lem:tameQuarticTrace}, followed by
the completed tensor operator of
Lemma~\ref{lem:ellipticCompletedTensor}.

Since $\sigma_{-1}(b)=-b$, opposite roots occur as differences.
Their norm divisors are
$\Gamma_{[g^{-1}]Q}-\Gamma_{-[g^{-1}]Q}$.
The Poincare computation in
Proposition~\ref{prop:ellipticCohenDetector} sends each such
difference to a common nonzero factor times
$\chi_0(g)^{-1}Q^*\eta$.  Keeping both elements of each opposite
pair, the full image is
\[
 \varepsilon\ell(\omega)
  \left(\sum_{g=1}^4\chi_0(g)^{-1}\sigma_g(b)\right)Q^*\eta
   =4\varepsilon\ell(\omega)e_{\chi_0}(b)Q^*\eta,
 \qquad \varepsilon\in\{1,-1\}.
\]
Here $\omega$ is the dual proper slope-one generator; the common
sign depends only on the Poincare and first-Chern-class conventions.
The factor of two for an opposite-root difference is already
accounted for by summing both elements of that pair.

All three factors are nonzero.  The first satisfies
$\ell(\omega)\ne0$ by construction, the second is nonzero by
Proposition~\ref{prop:dimtwoActualCoefficient}, and $Q^*\eta$
survives in the ordinary generic-Cohen de Rham quotient by
Lemma~\ref{lem:scalarFrobeniusCohen}.  A finite coefficient extension
preserves this nonzero vector.  Both operators send actual exact
forms to actual exact forms.  Thus $\Xi$ cannot have been exact.
This argument uses the exact character projection, so no
independence assertion for a tensor product of the entire root
Picard lattice with $\mathbf Q_5$ is needed.
\end{proof}

\begin{proof}[Proof of Theorem~\ref{thm:dimensiontwoMain}]
Proposition~\ref{prop:dimtwoActualCoefficient} supplies an actual
integrally antisymmetrized coefficient over
$W_{\rm ar}=\mathbf Z_{(5)}[\rho]$, with zero residue reduction.
Proposition~\ref{prop:dimtwoWeightedNonzero} verifies the ordinary
continuous de Rham nonvanishing required by
Proposition~\ref{prop:dimensiontwoactualassembly}.
That proposition constructs $c=i_*(\beta)$ in $K_3(A_2)$, proves that its
generic image is zero integrally, and proves infinite order for
both $c$ and its image in $K_3(\widehat A_2)$.
These are the two assertions of the theorem.
\end{proof}

\begin{remark}[The precise strength of the conclusion]
The residue field in this dimension-two construction is
$\mathbf F_5(T)$.  No dimension-two assertion with finite residue
field is made here.  Infinite order of the actual classes follows
from their nonzero images under homomorphisms to rationalized
completed spectra; injectivity of the map from actual $K$-theory
to its completion is neither asserted nor needed.
\end{remark}


% Independent optional supplement; proposed argument under internal review.
% Insert after the actual assembly and weighted detector in the dimension-two draft.
% Additional bibliography keys: Weibel, BunkeTamme, Borel, DeJeuWedge.
% The only use of the completed detector is the final paragraph.

\section{One explicit triple of units}
\label{sec:dimtwoExplicitTriple}

Retain the arithmetic local ring $A_2$ and its completion $\widehat A_2$
from \eqref{eq:dimensiontwoArithmeticSource}--\eqref{eq:dimensiontwoCompletedSource}.
The following supplement identifies an explicit class with a detected
weighted class \emph{after the specified test pullback}.

\begin{proposition}[Proposed explicit dimension-two class]
\label{prop:dimtwoExplicitTriple}
The triple
\begin{equation}\label{eq:dimtwoExplicitTriple}
 c_{\rm exp}=\left\{2+y,\ 1-x,\
 1+(1-x)(x^3+Ty^3)y\bigl((1+y)^5+1\bigr)\right\}\in K_3(A_2)
\end{equation}
has infinite order in $A_2$ and in $\widehat A_2$, and its
restriction to the localization at $x$ is zero integrally in both cases.
\end{proposition}

\begin{proof}
\emph{Integral vanishing.}
Put $t=-1-y$, $v=x^3+Ty^3$, and
\[
 H(t)=\frac{t^{10}-1}{\Phi_{10}(t)}=(t^5-1)(t+1)
      =y\bigl((1+y)^5+1\bigr),\quad b=vH(t).
\]
Then $xv=-\Phi_{10}(t)$, so $1-xb=t^{10}$.
Set $g=1-t$, $\lambda=1-x$, and $\mu=1+\lambda b$.
These are units: their residues are $2,1,1$ respectively, and
$t$ is also a unit. In the convention
$\langle r,s\rangle=\{r,1-rs\}$ when $r$ is a unit, the
Dennis--Stein relations give
\[
 \delta=\langle x,b\rangle
 =\langle x,b+1-xb\rangle=\langle x,\mu\rangle
 =\{\lambda t^{10},\mu\}.
\]
Here $\langle x,1\rangle=0$ and $1-x\mu=\lambda t^{10}$;
see \cite[Chapter~III, Definition~5.11 and Exercise~5.6]{Weibel}.
Since $\{g,t^{10}\}=10\{1-t,t\}=0$, multiplication gives
\begin{equation}\label{eq:dimtwoDSidentity}
 [g]\delta=\{g,\lambda,\mu\}=c_{\rm exp},\qquad
 [g]\delta|_{D(x)}=\{g,x,t^{10}\}=0.
\end{equation}
All identities are integral and persist under every ring map in use.

\emph{The universal supported coefficients.}
Write $E_0=\mathbf Q(\rho)$, where $\rho$ is primitive of order five,
and $t_g=-\rho^g$ for $g\in\mathbf F_5^\times$. On the smooth surface
\[
 Q^0=\operatorname{Spec}E_0[X,V,t]/(XV+\Phi_{10}(t)),\qquad
 U=Q^0[1/(t(1-t))],
\]
the class $\alpha=[1-t]\langle X,VH(t)\rangle$ is defined and
vanishes on $U[1/X]$. The entire divisor $D=V(X)$ is contained
in $U$. Localization and devissage therefore give an actual
supported lift $\kappa\in K_3(Q^0)$ of $\alpha$.
Put $\eta_g=1-[(X,t-t_g)]$. Since the four components of $D$
are affine lines, $\kappa=\sum_g\gamma_g\eta_g$ with
$\gamma_g\in K_3(E_0)$, and $\sum_g\eta_g=0$.

\emph{Independence of the supported lift.}
This lift is unique rationally as a class on $Q^0$.
Here is a direct $K$-theoretic check, independent of its periods,
using localization and the Laurent fundamental theorem \cite{Weibel}.
Put $Z=\mathbf A^1_{E_0}\setminus\{0,1\}$.  Then
$U[1/X]\simeq\mathbf G_m\times Z$, and localization on the
affine line, split by the units $t$ and $1-t$, gives
\[
 K_r(Z)\simeq K_r(E_0)\oplus
       K_{r-1}(E_0)[t]\oplus K_{r-1}(E_0)[1-t].
\]
Together with the Laurent fundamental theorem this gives
\[
\begin{split}
 K_4(U[1/X])\simeq{}&
 K_4(E_0)\oplus K_3(E_0)[t]\oplus K_3(E_0)[1-t]\\
 &{}\oplus K_3(E_0)[X]
 \oplus K_2(E_0)[X][t]\oplus K_2(E_0)[X][1-t].
\end{split}
\]
The first three summands extend to $U$, so their boundary in
$K_3(D)=K_3(E_0)^4$ is zero.  For a number field,
$K_2(E_0)_{\mathbf Q}=K_4(E_0)_{\mathbf Q}=0$: localization
reduces this to the rational even $K$-group vanishing for its
number ring and the torsion odd $K$-groups of finite fields.
The only remaining boundary is that of
$\beta[X]$, namely the diagonal tuple
$\pm(\beta,\beta,\beta,\beta)$ with a common convention-dependent
sign, since $X$ has order one on each root component.  Hence
\[
 \ker\bigl(K_3(D)_{\mathbf Q}\longrightarrow K_3(U)_{\mathbf Q}\bigr)
       =\Delta K_3(E_0)_{\mathbf Q}.
\]
The diagonal pushes forward to zero on $Q^0$, by the
length-one free resolution of the principal divisor $X=0$.
Consequently any two actual supported lifts of $\alpha$ give
the same class in $K_3(Q^0)_{\mathbf Q}$.
Moreover, since
$Q^0[1/X]\simeq\mathbf G_m\times\mathbf A^1_{E_0}$, the same
Laurent calculation gives
$\ker(K_3(D)\to K_3(Q^0))=\Delta K_3(E_0)$ already integrally.
Thus its supported coefficients are unique modulo the diagonal.

For a root $a$ of order ten define
$\beta_{\rm BT}(a)=\frac1{10}\operatorname{bl}(10[a])$
using \cite[Theorem~4.9 and (4.8)]{BunkeTamme}. Both $a$ and
$1-a$ are global units, and $10[a]$ has zero wedge boundary.
The normalization at every complex embedding $\sigma$ is
$r_2(\beta_{\rm BT}(a))_\sigma=-iD_{\rm BW}(\sigma a)$.
We claim the exact rational identity
\begin{equation}\label{eq:dimtwoUniversalCoefficients}
 \kappa=-10\sum_g\beta_{\rm BT}(t_g)\eta_g
          \quad\hbox{in }K_3(Q^0)_{\mathbf Q}.
\end{equation}

Here is the period calculation fixing all coefficients. For any two
roots at any embedding, their chord avoids zero, the other roots,
and $1$; its interior lies inside the unit circle. On its matching
sphere take $X=\sqrt{|\Phi_{10}(t)|}e^{i\theta}$ and
$V=-\Phi_{10}(t)/X$, adding the two poles. Multiplicativity and the
unit normalization of \cite[Proposition~2.8, Theorem~2.31 and
Lemma~4.5]{BunkeTamme} represent the regulator there by
\[
 -10\operatorname{Log}(1-t)\,\frac{dX}{X}\wedge\frac{dt}{t}.
\]
The form extends across the poles; equivalently its curvature is
$-dX\wedge d(VH)/(1-XVH)$. With the sphere oriented from root
$i$ to root $j$, its period is
$-20\pi i(\operatorname{Li}_2(\sigma t_j)
                -\operatorname{Li}_2(\sigma t_i))$ modulo
$\mathbf R(3)=i\mathbf R$. Its real part is
$20\pi(D_{\rm BW}(\sigma t_j)-D_{\rm BW}(\sigma t_i))$.
The periods of $\operatorname{ch}_1(\eta_i)$ and
$\operatorname{ch}_1(\eta_j)$ are $2\pi i$ and $-2\pi i$.
Thus, at every embedding,
\[
 r_2(\gamma_i-\gamma_j)
       =-10r_2\bigl(\beta_{\rm BT}(t_i)-\beta_{\rm BT}(t_j)\bigr).
\]
Borel injectivity \cite{Borel,BunkeTamme} proves the equality of
these rational $K_3(E_0)$ classes. Their remaining common
coefficient contributes zero since $\sum_g\eta_g=0$.
This proves \eqref{eq:dimtwoUniversalCoefficients}.

\emph{Comparison with the actual cyclotomic coefficient.}
The Besser--de Jeu symbol map and the BT normalization differ by
one scalar $s\in\mathbf Q^\times$, independent of the number
field and root. Indeed, \cite[Proposition~4.1 and Remark~5.2]{DeJeuWedge}
computes the former complex regulator at all embeddings as one
fixed nonzero real multiple of $iD_{\rm BW}$. Calibration at a
primitive sixth root in $\mathbf Q(\sqrt{-3})$, whose rational
$K_3$ has rank one, makes that multiple rational. Borel injectivity
then gives $\beta_{\rm BDJ}(a)=s\beta_{\rm BT}(a)$ at every root.
The field map is the same one used in
\cite[Theorem~1.10(1)--(2), Remark~1.7]{BesserDeJeu}.
Finite-residue localization identifies
$K_3(W_{\rm ar})_{\mathbf Q}$ with $K_3(E_0)_{\mathbf Q}$,
so this comparison also holds before local completion.
Consequently the coefficient and character argument of
Proposition~\ref{prop:dimtwoActualCoefficient} applies to
\[
 \Gamma=\operatorname{bl}(10[-\rho])\in K_3(\mathcal O_{E_0}),
 \qquad \beta=\Gamma-\sigma_{-1}\Gamma\in K_3(W_{\rm ar}).
\]
This is an actual integral antisymmetrization. Inversion of
$D_{\rm BW}$ at all embeddings and Borel injectivity show
$\beta_{\mathbf Q}=20\beta_{\rm BT}(-\rho)$ in $K_3(E_0)_{\mathbf Q}$.
In particular its two odd local character components are nonzero,
and \eqref{eq:dimtwoUniversalCoefficients} becomes
\begin{equation}\label{eq:dimtwoExactHalf}
 \kappa=-\frac12\sum_g\sigma_g(\beta)\eta_g.
\end{equation}

\emph{The test pullback.}
The substitutions $X=x$, $V=x^3+Ty^3$, $t=-1-y$ identify
$Q^0[1/(t+1)]$ with the generic surface used in the geometric
argument. For the completed argument they define a ring map
$\mathcal O(Q^0)\to B[1/\pi]$,
with $X\mapsto\pi X$, $t\mapsto-1-\pi Y$.
Both $t$ and $1-t$ are units already in $B$, and the root line
$(X,t-t_g)$ pulls back to $J_g[1/\pi]$.
The denominator and completion maps of
Proposition~\ref{prop:dimensiontwoactualassembly} therefore give
\[
 c_{\rm exp}|_{B[1/\pi]}
 =\widehat c_{\rm exp}|_{B[1/\pi]}
 =-\frac12\sum_g\sigma_g(\beta)(1-[J_g])\big|_{B[1/\pi]}.
\]
Proposition~\ref{prop:dimtwoWeightedNonzero} and the actual assembly
detect the right side as nonzero in rationalized completed
$K$-theory. Thus both actual source classes have infinite order.
This uses equality on the split surface and on the test algebra;
it asserts no equality between $c_{\rm exp}$ and a Gysin class
inside $K_3(A_2)$ itself.
\end{proof}

% Optional direct uncompleted verification: the chord from exp(-i*pi/5)
% to exp(i*pi/5) stays in Q^0[1/(t+1)] and gives positive real period
% 40*pi*D_BW(exp(i*pi/5)).  The proved arithmetic H^2 persistence then
% detects this class at every permitted finite localization, independently
% of the completed argument.

% Independent targeted extension, 7 October 2026.
% No main-manuscript edits. Requires the preceding quartic detector,
% its essentially smooth AMMN comparison, and the scalar Cohen lemma.

\begin{corollary}[The same explicit symbol in ordinary Milnor theory]
\label{cor:twodimexplicitMilnor}
For either $A=A_2$ or $A=\widehat A_2$, the ordinary Milnor symbol
\[
 m=\left\{2+y,\ 1-x,\
 1+(1-x)(x^3+Ty^3)y\bigl((1+y)^5+1\bigr)\right\}
       \in K_3^M(A)
\]
has infinite order and has zero image in
$K_3^M(\Frac A)$, integrally. In particular the ordinary
Milnor Gersten augmentation also fails after rationalization.
\end{corollary}

\begin{proof}
Write
\begin{align*}
 t&=-1-y,\qquad g=1-t,\qquad \lambda=1-x,\qquad q=t^{10},\\
 b&=(x^3+Ty^3)y\bigl((1+y)^5+1\bigr),\qquad \mu=1+\lambda b.
\end{align*}
The defining equation gives $1-xb=q$ and hence
$1-x\mu=\lambda q$. The elements $g,\lambda,\mu,q$
are units in $A$: their residues are respectively $2,1,1,1$.
In its fraction field, $x$ and $x\mu$ are also nonzero.
The two Steinberg relations and skew-symmetry in field
Milnor $K$-theory give
\begin{align*}
 0&=\{x\mu,\lambda q\},\qquad \{x,\lambda\}=0,\\
 \{x,q\}&=-\{\mu,\lambda\}-\{\mu,q\}
            =\{\lambda q,\mu\}.
\end{align*}
Furthermore $\{g,q\}=10\{1-t,t\}=0$. Therefore, without
rationalizing,
\[
 m=\{g,\lambda,\mu\}
   =\{g,\lambda q,\mu\}
   =\{g,x,q\}
   =-\{g,q,x\}=0
       \quad\text{in }K_3^M(\Frac A).
\]
The natural homomorphism $K_3^M(A)\to K_3(A)$ induced
by products of units sends $m$ to the already detected
infinite-order class $c_{\rm exp}$. Thus no nonzero integer
can kill $m$. This only uses the natural product homomorphism;
it does not use a comparison isomorphism or a retraction.
\end{proof}

\section{Visible coefficient rank and all higher odd degrees}
\label{sec:dimtwoVisibleHigherRank}

The elliptic detector sees one specified odd character.  Consequently
independence of unprojected anti-invariant coefficients is not by
itself sufficient.  We first construct arbitrary finite independent
families in the visible character.

\begin{lemma}[Actual coefficients of arbitrary visible rank]
\label{lem:dimtwoVisibleHigherCoefficientRank}
Fix $n\geq2$ and $m\geq2$. Put
\[
 V_m=W(\mathbf F_{5^m}),\quad F_m=V_m[1/5],\quad
 O_m=V_m[\rho],\quad E_m=F_m(\rho),\quad \pi=\rho-1,
\]
where $\rho$ is a primitive fifth root of unity.
Let $\sigma_g$ fix $F_m$ and send $\rho$ to $\rho^g$, and put
\[
 e_3=\frac14\sum_{g=1}^4\omega(g)^{-3}\sigma_g,\qquad
 \tau_3=\sum_{g=1}^4\omega(g)^{-3}\rho^g.
\]
There are $m-1$ actual classes
$\beta_j\in K_{2n-1}(O_m)$, each of the form
$(1-\sigma_{-1})\beta_{0,j}$, with the following properties:
their reductions in $K_{2n-1}(\mathbf F_{5^m})$ vanish integrally,
and their AMMN point coordinates $b_j=r_{n,O_m}(\beta_j)\in E_m$
have $\mathbf Q_5$-linearly independent projections $e_3b_j$.
\end{lemma}

\begin{proof}
Define
$S_j(z)=(z\,d/dz)^j(1-z)^{-1}\in\mathbf Z[z,(1-z)^{-1}]$.
The same expansion and continuation argument used in
Lemma~\ref{lem:dimtwoDilogCharacters}, now with
$\binom{n+j-1}{j}$ in place of $j+1$, gives
\[
 F_n(z):=\operatorname{Li}_n(z)-5^{-n}\operatorname{Li}_n(z^5)
 =\sum_{j\geq0}(-1)^j\binom{n+j-1}{j}5^j S_j(z^5)
             \sum_{a=1}^4\frac{z^a}{a^{n+j}}.
\]
The series is uniformly convergent on $|z|\leq1$,
$|1-z^5|=1$, and agrees there with Coleman continuation by
\cite[Theorem~2.7(3)]{BesserDeJeuAlgorithm}.
This again uses power substitution, not a Frobenius automorphism
of the ramified field.

For a Teichmuller root
$\zeta\in\mu_{5^m-1}(V_m)\setminus\{1\}$, the value
$(\zeta\rho)^5=\zeta^5$ is fixed by every $\sigma_g$.
The exact Gauss identity from the preceding coefficient lemma
therefore gives
\begin{equation}
\label{eq:dimtwoVisibleHigherDepletion}
\begin{split}
 U_n(\zeta)
 &:=\frac{4e_3\operatorname{Li}_n(\zeta\rho)}{\tau_3}\\
 &=\sum_{j\geq0}(-1)^j\binom{n+j-1}{j}5^j S_j(\zeta^5)
        \sum_{a=1}^4\frac{\zeta^a\omega(a)^3}{a^{n+j}}
       \in V_m.
\end{split}
\end{equation}
The terms with $j\geq1$ belong to $5V_m$.  Thus, writing
$t=\overline\zeta$, its reduction is
\[
 h_{3,n}(t)
   =\frac{P_n(t)}{1-t^5},\qquad
 P_n(t)=\sum_{a=1}^4a^{3-n}t^a\in\mathbf F_5[t],
 \qquad t\in\mathbf F_{5^m}^{\times}\setminus\{1\}.
\]
The polynomial $P_n$ is nonzero, has coefficient one at $t$,
and has degree four.  Hence $h_{3,n}$ is nonconstant and has
rational-map degree at most five.  Every finite fiber has at most
five elements, so
\[
 \#h_{3,n}\bigl(\mathbf F_{5^m}^{\times}\setminus\{1\}\bigr)
 \geq\left\lceil\frac{5^m-2}{5}\right\rceil=5^{m-1}.
\]
Its $\mathbf F_5$-span has dimension at least $m-1$.
Choose $m-1$ independent image values and corresponding
Teichmuller roots $\zeta_j$.  Reduction of a normalized
$\mathbf Q_5$-linear relation shows that the $U_n(\zeta_j)$
are $\mathbf Q_5$-independent.  Since $\tau_3\ne0$, the values
$e_3\operatorname{Li}_n(\zeta_j\rho)$ are independent as well.

These are regulators of actual rational number-field classes.
Indeed, all $\zeta_j\rho$ lie in
$L=\mathbf Q(\mu_{5^m-1},\rho)$ and are roots of unity.
At the chosen prime above $5$ they and their complements are
units, since $\overline\zeta_j\ne1$.  The rational symbols
$[\zeta_j\rho]_n$ are cycles, and
\cite[Theorem~1.10(2)]{BesserDeJeu} maps them to a common
nonzero rational multiple of
$L_{\mathrm{mod},n}(\zeta_j\rho)
 =\operatorname{Li}_n(\zeta_j\rho)$.
There is no Beilinson--Soule assumption for this number field.
The local completion at the chosen prime is $E_m$.

Localization at that coefficient prime lifts these rational
classes to its integer ring, since
$K_{2n-2}(\mathbf F_{5^m})=0$.
Clear one common integer denominator and then form the actual
differences
$\beta_j=(1-\sigma_{-1})\beta_{0,j}$, mapping them to $K_{2n-1}(O_m)$.
Their residue reductions are zero because $\sigma_{-1}$ fixes
the residue field.  Applying $e_3$ to their syntomic regulators
multiplies each of the preceding values by the same nonzero
rational scalar, including the factor two from the difference.

The ramified point comparison
\cite[Example~2.2.4 and Propositions~2.2.7, 2.2.9, 2.3.4]{HuberKings}
is Galois-equivariant and identifies syntomic and etale point
cohomology in every twist $n\geq2$.  Thus the $e_3$-projections
of the etale Chern values are $\mathbf Q_5$-independent.
The finite-coefficient construction, exactly as in
Proposition~\ref{prop:dimtwoActualCoefficient}, gives a natural
Galois-equivariant $\mathbf Q_5$-linear map
\[
 \pi_{2n-1}\bigl(K(O_m)^\wedge_5[1/5]\bigr)
       \longrightarrow H^1(E_m,\mathbf Q_5(n))
\]
through which the actual Chern class factors.
For clarity, this follows from
\cite[Proposition~2.1.1, Lemma~2.3 and Propositions~2.4, 2.8]{WeibelChern}:
the finite Chern maps are additive, natural, and compatible with
coefficient reduction.  Their target inverse limit is continuous
cohomology because the finite $H^0$ groups are Mittag--Leffler.
Only the canonical source projections are used.

The point case of the AMMN comparison identifies the source with
$t^{1-n}E_m$, naturally under $\sigma_g$.  A relation among the
$e_3b_j$ would therefore induce the same relation among the
independent projected Chern values.  This proves the claim
without identifying the AMMN and polylogarithm coordinates.
\end{proof}

\begin{remark}
The rational-map degree bound cannot uniformly be replaced by four.
When $n\equiv3\pmod4$ one has
\[
 h_{3,n}(t)=(t-1)^{-5}-(t-1)^{-1}.
\]
Its image lies in the trace-zero hyperplane of
$\mathbf F_{5^m}/\mathbf F_5$.  The lower bound $m-1$ accommodates
this actual loss of one dimension.
\end{remark}

\begin{theorem}[One complete quartic ring in every odd degree]
\label{thm:dimtwoCompleteAllOddInfiniteRank}
Put
\[
 k=\overline{\mathbf F}_5,\qquad W=W(k),\qquad
 V=\widehat{W[T]_{(5)}}^{\,(5)},\qquad
 R=V[[x,y]]/(\Phi_5(1+y)+x^4+Txy^3).
\]
Then $R$ is a complete regular local domain of dimension two,
with maximal ideal $(x,y)$ and residue field $k(T)$.
For every $n\geq2$,
\[
 \dim_{\mathbf Q}\ker\left(
 K_{2n-1}(R)_{\mathbf Q}\longrightarrow
 K_{2n-1}(\operatorname{Frac}R)_{\mathbf Q}\right)=\infty.
\]
For each fixed $n,m\geq2$, the kernel contains $m-1$ independent
classes represented by actual integral classes that vanish on
$R[1/x]$ integrally.  The ring $R$ is independent of $n,m$.
\end{theorem}

\begin{proof}
The regularity argument is the same linear-term calculation as in
Proposition~\ref{prop:dimensiontwoactualassembly}.
Set $W'=W[\rho]$ and $E_\infty=W'[1/5]$.
For fixed $n,m$, take the actual coefficients of
Lemma~\ref{lem:dimtwoVisibleHigherCoefficientRank}.
Naturality of the AMMN point calculation sends their coordinates
under $E_m\hookrightarrow E_\infty$; hence their visible
projections remain $\mathbf Q_5$-independent.
No filtered-colimit description of completed $K$-theory is needed.

Use the same finite parameter cover and scaled formal test ring
$B$ as in \eqref{eq:dimensiontwotestalgebra}, now over $V$.
It is the completion of an essentially smooth $W'$-algebra
of relative dimension two.  The spectrum comparison underlying
Lemma~\ref{lem:dimensiontwoAMMNinput} gives, for each fixed $n$,
the relative character to $HC_{2n-2}^{\mathrm{cts}}(B/W';\mathbf Q_5)$.
In its finite mixed-HKR model the top-form summand is
\begin{equation}
\label{eq:dimtwoHigherTopColumn}
 t^{2-n}\left(
 \Omega^2_{B/W'}[1/5]/d\Omega^1_{B/W'}[1/5]\right).
\end{equation}
Indeed its cycles in this summand are all two-forms, and its
boundaries are precisely $t^{2-n}d\Omega^1$; form degrees above
two vanish.

The zero residue reductions of the coefficients provide actual
relative lifts at the point.  Pull them to $B$ and multiply by
$\eta_g=1-[J_g]$.  The same negative-cyclic module product gives
\[
 (t^{1-n}\sigma_g(b_j))\,(t\gamma_g)
           =t^{2-n}\sigma_g(b_j)\gamma_g.
\]
Rank zero removes the scalar term, higher Chern forms vanish in
dimension two, and the cyclic-shuffle terms vanish because the
point coefficient is a scalar relative to $W'$.
This argument also applies when the initial cyclic column is
negative; it does not discard positive columns without a degree
or scalar-bar calculation.

The trace and contraction calculation in
Proposition~\ref{prop:dimtwoWeightedNonzero} now detects the
coefficient of \eqref{eq:dimtwoHigherTopColumn} as
\[
 4\varepsilon\ell(\omega)e_3(b_j)\,Q^*\eta.
\]
The fixed vector $Q^*\eta$ is nonzero over $E_\infty$ after the
generic Cohen completion, by
Lemma~\ref{lem:scalarFrobeniusCohen}.
Thus these character images are $\mathbf Q_5$-independent.

Write $C=B/\pi$.  As before, $C$ is a filtered colimit of smooth
affine surfaces over $k$.  By
\cite[Theorem~8.4]{GeisserLevineCharP},
$\pi_q\mathcal K(C)=0$ for every $q\geq3$.
The relative and regular-divisor localization sequences therefore
give, in the present degree,
\[
 \pi_{2n-1}\mathcal K(B,(\pi))
 \hookrightarrow\pi_{2n-1}\mathcal K(B)
 \hookrightarrow\pi_{2n-1}\mathcal K(B[1/\pi]).
\]
These use the vanishings in degrees $2n$ and $2n-1$, respectively.
No vanishing assertion about actual $K_{2n}(C)$ or
$K_{2n-1}(C)$ is needed.

The source divisor $x=0$ has ring $V[\rho]$.
Pull each $\beta_j$ to that divisor and push forward to define
$c_j\in K_{2n-1}(R)$.  Localization gives $c_j|_{R[1/x]}=0$
integrally.  Finite flat coefficient extension $V\to V[\rho]$,
followed by inversion of $5$, splits the divisor into the four
root components with coefficients $\sigma_g(\beta_j)$.
Pullback of their line-bundle classes under the scaled map gives
\[
 c_j|_{B[1/\pi]}
   =\sum_g\sigma_g(\beta_j)(1-[J_g])\big|_{B[1/\pi]}.
\]
This is exactly the actual product detected above; the scaled map
itself is not required to be flat.
Consequently the $c_j$ are rationally independent.
Since $m$ is arbitrary while $R$ is fixed, the kernel has infinite
dimension for this $n$.  Applying the argument separately to each
$n\geq2$ proves the theorem.
\end{proof}

% Independent insertion, 7 October 2026.
% Requires the visible coefficient rank lemma and the previously constructed
% one-parameter quartic detector. No main-manuscript edits.

\section{One complete two-dimensional ring in every degree at least three}
\label{sec:dimtwoAllDegreesFixedRing}

The additional parameter below is a unit of the source ring.
Its logarithmic differential raises the degree of the detected
class by one.  Detecting the even products also proves independence
of their odd factors, including degree three.

\begin{lemma}[Residue after full generic Cohen completion]
\label{lem:dimtwoFullCohenLaurentResidue}
Let $W'$ be the coefficient discrete valuation ring, with
uniformizer $\pi$, used in the quartic construction.  Let $D_0$
be a smooth affine $W'$-curve whose special fiber is integral,
and put
\[
 A_T=\widehat{(D_0)_{(\pi)}}^{\,(\pi)},\qquad
 A_{TS}=\widehat{(D_0[S])_{(\pi)}}^{\,(\pi)}.
\]
Here localization at $(\pi)$ inverts every element with nonzero
reduction.  Let $B_T$ be the $\pi$-adic completion of a smooth
affine relative curve over $A_T$, defined by a model over $D_0$,
and let $B_{TS}$ be the corresponding completed model over
$A_{TS}$.  Suppose their continuous differentials over $W'$
are finite projective of ranks two and three, respectively.
There is a $W'$-linear operator
\[
 \operatorname{Res}_S:
 \Omega^{r,\mathrm{cts}}_{B_{TS}/W'}
       \longrightarrow \Omega^{r-1,\mathrm{cts}}_{B_T/W'}
\]
which is integral and $\pi$-adically continuous and satisfies
\[
 \operatorname{Res}_S(d\alpha)=-d\operatorname{Res}_S(\alpha),
 \qquad
 \operatorname{Res}_S(\theta\wedge d\log S)=\theta
       \quad\text{for every }\theta\in\Omega^2_{B_T/W'}.
\]
In particular it sends actual exact three-forms to actual exact
two-forms.  Thus
\[
 \frac{\Omega^2_{B_T/W'}[1/5]}
      {d\Omega^1_{B_T/W'}[1/5]}
 \longrightarrow
 \frac{\Omega^3_{B_{TS}/W'}[1/5]}
      {d\Omega^2_{B_{TS}/W'}[1/5]},
 \qquad [\theta]\longmapsto[\theta\wedge d\log S],
\]
is injective; these are ordinary continuous quotients.
\end{lemma}

\begin{proof}
For every $a\geq1$ there is a compatible ring map
\[
 B_{TS}/\pi^a
       \longrightarrow (B_T/\pi^a)((S)).
\]
Indeed, if $g\in D_0[S]$ has nonzero reduction, its image modulo
$\pi$ is a nonzero polynomial over the field
$\operatorname{Frac}(D_0/\pi)$.  It is therefore a unit in that
field's Laurent series ring.  The ideal generated by $\pi$ in
$(B_T/\pi^a)((S))$ is nilpotent, so this inverse lifts.  This
argument applies to every primitive denominator, including
denominators whose literal lowest coefficient is not a unit.
The lifts of the inverses are unique and commute with reduction
in $a$.  The maps extend to the completed algebra because the
inverse-limit target is complete.  Thus the relevant target is
\[
 \mathcal L=\varprojlim_a(B_T/\pi^a)((S)).
\]
No map into the ordinary ring $B_T((S))$ is asserted.

Use Laurent series of continuous base forms, with differential
defined coefficientwise and with the usual $S$-derivative.
These operations are compatible modulo $\pi^a$ and preserve
the $\pi$-adic filtration.  The universal continuous derivation
of $B_{TS}$ therefore gives a map into this Laurent differential
complex.  Finite projectivity of the source form modules
justifies extension to their completions and exterior powers.

Write a form of degree $r$ in the target as
$\alpha(S)+\beta(S)\wedge dS$, with base-form degrees $r$
and $r-1$.  Define
\[
 \operatorname{Res}_{S,r}
       \bigl(\alpha(S)+\beta(S)\wedge dS\bigr)
      =(-1)^{r-1}[S^{-1}]\beta(S).
\]
Coefficient extraction is compatible with reduction and is
integral.  The coefficient of $S^{-1}$ in the derivative of a
Laurent series is zero.  It follows directly that
$\operatorname{Res}_{S,r+1}d=-d\operatorname{Res}_{S,r}$.
For $r=3$ its sign is positive, which gives the second displayed
identity.  Inverting $5$ preserves these identities.  All
quotients in the conclusion use the ordinary image of the
differential.
\end{proof}

\begin{theorem}[A fixed complete ring in all degrees]
\label{thm:dimtwoCompleteAllDegreesInfiniteRank}
Put
\[
 k=\overline{\mathbf F}_5,\qquad W_0=W(k),\qquad
 V_{TS}=\widehat{W_0[T,S]_{(5)}}^{\,(5)},
\]
and
\[
 R_{TS}=V_{TS}[[x,y]]/
       \bigl(\Phi_5(1+y)+x^4+Txy^3\bigr).
\]
Then $R_{TS}$ is a complete regular local domain of dimension
two, with maximal ideal $(x,y)$ and residue field $k(T,S)$.
For every integer $j\geq3$,
\[
 \dim_{\mathbf Q}\ker\!\left(
 K_j(R_{TS})_{\mathbf Q}\longrightarrow
 K_j(\operatorname{Frac}R_{TS})_{\mathbf Q}\right)=\infty.
\]
For every $n,m\geq2$, the odd degree $2n-1$ and the even degree
$2n$ kernels contain $m-1$ independent classes represented by
actual integral classes which vanish on $R_{TS}[1/x]$ integrally.
The ring $R_{TS}$ is independent of $n$ and $m$.
\end{theorem}

\begin{proof}
The ring $W_0[T,S]_{(5)}$ is a discrete valuation ring with
uniformizer $5$ and residue field $k(T,S)$.  Its completion
$V_{TS}$ has the same properties.  The defining equation of
$R_{TS}$ is congruent to $5$ modulo $(5,x,y)^2$ in the regular
local ambient ring of dimension three.  Its quotient is
therefore regular of dimension two, and is complete with
maximal ideal $(x,y)$.  In particular it is a domain.
The element $S$ is a unit of $R_{TS}$.

Set $W'=W_0[\rho]$, $\pi=\rho-1$, and
$E_\infty=W'[1/5]$.  Fix $n,m\geq2$, and take the $m-1$
actual coefficients $\beta_j$ from
Lemma~\ref{lem:dimtwoVisibleHigherCoefficientRank}.
Map them to $K_{2n-1}(W')$ and write $b_j\in E_\infty$
for their natural AMMN point coordinates.  Their visible
projections $e_3b_j$ remain $\mathbf Q_5$-independent.
Their reductions vanish integrally, since each coefficient
is an actual difference $(1-\sigma_{-1})\beta_{0,j}$.

Choose the same finite parameter cover, independent of $S$,
as in the one-parameter quartic detector.  In the notation of
Lemma~\ref{lem:dimtwoFullCohenLaurentResidue}, it gives the old
coefficient ring $A_T$ and the new one $A_{TS}$.  The latter
receives $V_{TS}[\rho]$: every primitive polynomial in $T,S$
has nonzero image in $k(U)(S)$, and hence is a unit.  Completion
extends this map.  Put
\[
 B_{TS}=A_{TS}\langle X,Y\rangle/
 \left(X^4+\frac{\Phi_5(1+\pi Y)}{\pi^4}+TXY^3\right),
 \qquad C_{TS}=B_{TS}/\pi.
\]
The old algebra $B_T$ is given by the same formula over $A_T$.
They are the completions of essentially smooth $W'$-algebras
of relative dimensions three and two, respectively.
Smoothness holds near the entire special fiber; localizing a
model at an element congruent to one modulo $\pi$ leaves the
completion unchanged.  In particular $B_{TS}$ and $C_{TS}$
are regular and $\pi$ is a non-zero-divisor.

For $g\in\mathbf F_5^\times$, let
\[
 a_g=\frac{\rho^g-1}{\pi},\qquad
 J_g=(X,Y-a_g),\qquad
 \eta_g=1-[J_g],\qquad \gamma_g=-c_1(J_g).
\]
These invertible ideals are pulled from $B_T$.  We can therefore
represent their Chern classes by two-forms pulled from the old
continuous complex, independently of $S$.

The actual residue sequence at the point supplies actual
relative lifts of all $\beta_j$; here
$K_{2n}(k)=0$ by Quillen's finite-field computation and filtered
colimits.  Pull these lifts to $B_{TS}$ and form their relative
products
\[
 \xi_j=\sum_g\sigma_g(\beta_j)\eta_g[S]
       \quad\text{in }K_{2n}(B_{TS},(\pi)),
\]
where the coefficient symbols in this formula denote the
specified relative lifts.

Apply the essentially smooth relative fiber comparison of
Lemma~\ref{lem:essentiallysmoothAMMN} in each fixed degree.
Its $K(B_{TS})$-module compatibility gives the negative-cyclic
action on the ordinary cyclic target.  In relative dimension
three the normalized finite mixed-HKR map uses only the
denominators $2$ and $6$, both units at $5$.
The point coefficient is $t^{1-n}\sigma_g(b_j)$, the rank-zero
line character starts with $t\gamma_g$, and the unit character
starts with $d\log S$.  Hence
\[
 \chi_{B_{TS}}(\xi_j)_{\mathrm{top}}
    =t^{2-n}\sum_g\sigma_g(b_j)\gamma_g\wedge d\log S
       \quad\text{in }HC^{\mathrm{cts}}_{2n-1}(B_{TS}/W';
                                                    \mathbf Q_5).
\]
Here $|t|=-2$.  This formula also holds when the initial cyclic
column is negative.  To check it at chain level, the rank-zero
line has underlying Hochschild degrees at least two; the
lowest unit term has degree one.  Every higher line or unit
term, or cyclic-shuffle correction in their product, has final
Hochschild degree at least five, and is killed by the
dimension-three HKR map.  The scalar point cup has no
cyclic-shuffle correction because its scalar lies in the
relative base.  Thus precisely the displayed ordinary
shuffle survives.

The top-form summand is
\[
 t^{2-n}\bigl(
 \Omega^3_{B_{TS}/W'}[1/5]/
                  d\Omega^2_{B_{TS}/W'}[1/5]\bigr).
\]
All its three-forms are cycles, and its boundaries are precisely
the indicated ordinary exact forms.  Apply
Lemma~\ref{lem:dimtwoFullCohenLaurentResidue} to their coefficients.
It recovers
$\sum_g\sigma_g(b_j)\gamma_g$ on $B_T$.
The old weighted trace and elliptic contraction of
Proposition~\ref{prop:dimtwoWeightedNonzero} sends these to
\[
 4\varepsilon\,\ell(\omega)e_3(b_j)\,Q^*\eta.
\]
The fixed base vector $Q^*\eta$ survives the old generic Cohen
completion by Lemma~\ref{lem:scalarFrobeniusCohen}, and remains
nonzero after the finite constant extension adjoining $\rho$.
The other fixed scalar $\ell(\omega)$ is nonzero.
Consequently the characters of the $\xi_j$ are
$\mathbf Q_5$-independent.

The special fiber $C_{TS}$ is a filtered colimit of smooth
affine threefolds over the perfect field $k$.  Geisser--Levine
\cite[Theorem~8.4]{GeisserLevineCharP} gives
\[
 K_q(C_{TS};\mathbf Z/5^\nu)=0
      \quad(q>3,\ \nu\geq1).
\]
The Milnor sequence for spectral completion, using degrees
$q$ and $q+1$, therefore gives
$\pi_q\mathcal K(C_{TS})=0$ for $q\geq4$.
For the even degree under consideration the relative and
regular-divisor localization sequences give
\[
 \pi_{2n}\mathcal K(B_{TS},(\pi))
   \hookrightarrow\pi_{2n}\mathcal K(B_{TS})
   \hookrightarrow\pi_{2n}\mathcal K(B_{TS}[1/\pi]).
\]
The first injection uses the vanishing in degree $2n+1$,
and the second that in degree $2n$; both are valid for $n\geq2$.
Here $\mathcal K=K^\wedge_5[1/5]$, as before.  No actual
special-fiber $K$-group vanishing is required.

It remains to match these detected even classes with actual
classes of the fixed source.  Its regular Cartier divisor
$i:x=0\hookrightarrow\operatorname{Spec}R_{TS}$ has ring
$V_{TS}[\rho]$.  Pull the coefficients to this divisor and set
\[
 c_{n,j}=i_*(\beta_j)\in K_{2n-1}(R_{TS}),\qquad
 a_{n,j}=c_{n,j}[S]\in K_{2n}(R_{TS}).
\]
Both classes vanish on $R_{TS}[1/x]$ integrally.
Make the finite flat coefficient extension
$V_{TS}\to V_{TS}[\rho]$, and then invert $5$.
The divisor $x=0$ splits into its four root components with
coefficients $\sigma_g(\beta_j)$.
Their length-one Cartier resolutions and the projection
formula express the resulting odd class as
$\sum_g\sigma_g(\beta_j)(1-[I_g])$, where
$I_g=(x,y-(\rho^g-1))$.
Now use the continuous scaled map
\[
 x\longmapsto\pi X,\qquad y\longmapsto\pi Y,\qquad
 R_{TS}\longrightarrow B_{TS}.
\]
It exists because the defining equation maps to $\pi^4$ times
the test equation and the maximal ideal maps into $(\pi)$.
After inverting $\pi$, pullback of the invertible ideals $I_g$
gives $J_g$.  Therefore
\[
 a_{n,j}|_{B_{TS}[1/\pi]}
   =\sum_g\sigma_g(\beta_j)\eta_g[S]\big|_{B_{TS}[1/\pi]}
       \quad\text{in actual }K_{2n}.
\]
Only the initial coefficient extension is asserted to be flat;
the subsequent substitution pulls back line bundles.
This equality and the two completed injections prove that
the actual classes $a_{n,j}$ are rationally independent.

Finally, a rational relation among the actual odd classes
$c_{n,j}$ would, after multiplication by the same actual unit
class $[S]$ in the same ring, give a relation among the
$a_{n,j}$.  Thus the odd classes are independent as well.
This argument includes $n=2$ and requires no vanishing of
$\pi_3\mathcal K(C_{TS})$.
Letting $m$ grow proves infinite rank separately in degrees
$2n-1$ and $2n$.  These pairs cover every $j\geq3$.
\end{proof}

\begin{remark}
No general injectivity of multiplication by $[S]$ is used:
it is proved only on the displayed finite spans by detecting
their products.  In particular the odd degree-three conclusion
does not use a degree-three generic-localization injection for
the new special fiber.  The residue field in this fixed-ring
statement is $\overline{\mathbf F}_5(T,S)$.
\end{remark}

\appendix
\section{The essentially smooth relative fiber comparison}
\label{sec:dimtwoAMMNAppendix}
% Independent insertion. Main manuscripts unchanged.
\begin{lemma}[The relative fiber comparison for an essentially smooth model]
\label{lem:essentiallysmoothAMMN}
Let $W_0=W(\overline{\mathbf F}_p)$ and let $W$
be a finite totally ramified extension, with uniformizer $\pi$
and fraction field $E$. Let $B_0$ be a localization of a smooth
$W$-algebra, with $\Omega^1_{B_0/W}$ finite projective of constant
rank $D$, and put $B=\widehat{B_0}_p$.
There is a natural actual relative character
\[
 \chi_B:K(B,\pi B)\longrightarrow
                 \Sigma HC^{\mathrm{cts}}(B/W;\mathbf Q_p),
\]
which factors through an equivalence from
$K(B,\pi B)^{\wedge}_p[1/p]$. It is $K(B)$-linear through the negative-cyclic
character, and the target has the finite continuous mixed-HKR
model, with ordinary de Rham quotients.
\end{lemma}

\begin{proof}
We justify the extension beyond topological finite presentation.
For each $r$, the ring $W_0/p^r$ is the filtered union of the
finite \emph{\'{e}tale} $\mathbf Z/p^r$-algebras
$W_r(\mathbf F_{p^m})$. Flat base change for cotangent complexes
therefore gives $\widehat L_{W_0/\mathbf Z,p}=0$.
Write $W=W_0[\pi]$ and let $f$ be the Eisenstein polynomial of
$\pi$. Then
\[
 \widehat L_{W/\mathbf Z,p}
   \simeq[W\xrightarrow{f'(\pi)}W]
\]
in cohomological degrees $-1,0$. Choose $a$ with
$p^a\in f'(\pi)W$; multiplication by $p^a$ on this complex is
null-homotopic.

Cotangent transitivity, after derived $p$-completion, gives
\[
 M\longrightarrow\widehat L_{B_0/\mathbf Z,p}
       \longrightarrow P,
 \qquad M=[B\xrightarrow{f'(\pi)}B],\quad
 P=\widehat\Omega^1_{B_0/W}.
\]
The target $P$ is finite projective in degree zero, so this
triangle splits as a triangle of underlying $B$-module
complexes. The error in the $i$th derived exterior power of its
projection is a sum of terms
$(\mathbf L\Lambda^rM)\otimes\Lambda^{i-r}P$, $1\le r\le i$.
The $r$th such term is killed by $p^{ar}$ as an object: derived
exterior power sends the null-homotopy of $p^a\operatorname{id}_M$
to one for $p^{ar}$.

The connective HKR filtration on Hochschild homology has graded
terms $(\mathbf L\Lambda^iL)[i]$. Any fixed homotopy range
involves only finitely many of these terms, even after completion.
Indeed derived $p$-completion preserves connectivity: for a
$c$-connective spectrum $X$, its only possible new bottom group
is $\lim^1_r(\pi_cX/p^r)=0$, because the quotient tower is
surjective. Thus the increasingly connective HKR tails remain
increasingly connective. The preceding bounds
and a finite extension argument therefore make
\[
 HH(B_0;\mathbf Z_p)\longrightarrow HH(B_0/W;\mathbf Z_p)
\]
a quasi-isogeny in each finite range. Its rationalization is an
equivalence. The bounds use the base different and the range;
they are unaffected by the localization set. No interchange of
an infinite localization colimit with completion is used.

Circle orbits on connective spectra involve only finitely many
cells in each fixed homotopy range. More explicitly, a finite
bar skeleton is built from finite colimits of
$(S^1)^j_+\wedge X$, which commute with completion since
$(S^1)^j$ is a finite CW complex. The omitted tail has arbitrarily
large connectivity, still so after completion. As in the proof of
\cite[Proposition~4.10]{AMMN}, this gives an equivalence of the
corresponding $HC(-;\mathbf Q_p)$ terms. Hence the square from
absolute to relative $HC^-\to HP$ is cartesian, by the norm
fiber sequence. This argument needs no assertion that its two
vertical arrows are individually equivalences.

By \cite[Proposition~4.2]{AMMN}, including the relative variant
stated just before Proposition~4.10, both $B_0$ and $B$ compute
the same continuous $p$-completed HH and cyclic theories.
The ordinary ring $B$ is henselian along $(p)$, so
\cite[Theorem~A]{AMMN} applies. The map $B/p\to B/\pi$ has
nilpotent kernel, and \cite[Corollary~3.8]{AMMN} identifies its
K-theory and TC after rationalized $p$-completion. Composing the
actual-to-completed map on relative K-theory with the resulting
fiber equivalence gives $\chi_B$.

The module argument is unchanged: for $H=THH(B)$, the objects
$H\otimes D$ and maps $\operatorname{id}_H\otimes f$ in
\cite[Proposition~2.8, Corollary~2.10, and Theorem~2.12]{AMMN}
form a diagram of $H$-modules in cyclotomic spectra.
The lax monoidal TC and circle functors give module diagrams,
their equivalences can be inverted there, and the relative HH
projection preserves the action. Together with the
multiplicative cyclotomic trace \cite[Theorem~7.4]{BGTTrace}, this gives the asserted actual
$K(B)$-linearity. This construction has no finite presentation
requirement.

Finally, the integral mixed HKR map
\cite[Definition~2.1, Remark~2.2 and Proposition~2.3]{BokstedtOttosen}
in degree $j\le D$ has
coefficient $D!/j!$ and induces multiplication by $D!$ on
$HH_j(B_0/W)=\Omega^j_{B_0/W}$. Its equivariant cone has a
Postnikov tower of length $D+1$ whose graded objects are killed
by $p^{v_p(D!)}$. These heart objects have trivial circle action;
the extension argument kills the entire equivariant cone by
$p^{(D+1)v_p(D!)}$. Thus circle orbits, fixed points, Tate
constructions, and derived completion retain only bounded
torsion error. The rational mixed-HKR comparison follows.
Each total degree has finitely many flat form modules, with
surjective reduction towers, so completion is termwise and
cohomology uses the ordinary image of $d$.
\end{proof}

\clearpage
{\small
\begin{thebibliography}{99}
\bibitem{DavisLangerZink}
C.~Davis, A.~Langer, and T.~Zink, \emph{Overconvergent de Rham--Witt
cohomology}, Ann. Sci. \'Ecole Norm. Sup. (4) \textbf{44} (2011),
197--262.
\href{https://doi.org/10.24033/asens.2143}{doi:10.24033/asens.2143}.

\bibitem{BesserDeJeuAlgorithm}
A.~Besser and R.~de~Jeu, \emph{Li$^{(p)}$-service? An algorithm for computing $p$-adic polylogarithms}, Math. Comp. \textbf{77} (2008), 1105--1134. \href{https://doi.org/10.1090/S0025-5718-07-02027-3}{doi:10.1090/S0025-5718-07-02027-3}.

\bibitem{BesserDeJeu}
A.~Besser and R.~de Jeu, \emph{The syntomic regulator for the $K$-theory of fields}, Ann. Sci. \'Ecole Norm. Sup. (4) \textbf{36} (2003), 867--924. Equation~(1.2) and Theorems~1.6, 1.10, 1.12.
\href{https://www.numdam.org/article/ASENS_2003_4_36_6_867_0.pdf}{Published article}.

\bibitem{HuberKings}
A.~Huber and G.~Kings, \emph{A $p$-adic analogue of the Borel regulator and the Bloch--Kato exponential map}. Section~2 gives the Chern-compatible comparison with Besser's syntomic cohomology and the exponential, without an unramifiedness restriction.
\href{https://arxiv.org/abs/math/0612611}{arXiv:math/0612611}.

\bibitem{WeibelChern}
C.~Weibel, \emph{\'Etale Chern classes at the prime 2}. Theorem~3.11(i) gives the line-bundle product formula in $K$-degree three; Section~5 recalls the indecomposable $K_3$ comparison.
\href{https://sites.math.rutgers.edu/~weibel/archive/papers-dir/chernclass.pdf}{Author's paper}.

\bibitem{SilvermanAEC}
J.~H.~Silverman, \emph{The Arithmetic of Elliptic Curves}, second edition,
Graduate Texts in Mathematics \textbf{106}, Springer, 2009;
Theorem~V.4.1(a).
\href{https://doi.org/10.1007/978-0-387-09494-6}{doi:10.1007/978-0-387-09494-6}.

\bibitem{CaroDAddezioInjectivity}
D.~Caro and M.~D'Addezio, \emph{Injectivity failure in crystalline comparisons},
arXiv:2511.11444; author's manuscript consulted 7 October 2026.
Theorems~3.2.1, 7.2.1, 7.3.1 and~7.3.6.
\href{https://carod.users.lmno.cnrs.fr/Injectivity.pdf}{Author's manuscript}.

\bibitem{AMMN}
B.~Antieau, A.~Mathew, M.~Morrow, and T.~Nikolaus,
\emph{On the Beilinson fiber square}. Propositions~2.8, 4.10, 4.12,
Corollary~2.10, Theorem~2.12, and Constructions~4.6, 4.11.
\href{https://arxiv.org/abs/2003.12541}{arXiv:2003.12541}.

\bibitem{GeisserLevineCharP}
T.~Geisser and M.~Levine, \emph{The $K$-theory of fields in
characteristic $p$}, Inventiones Mathematicae \textbf{139}
(2000), 459--493. Theorem~8.4.
\href{https://www2.rikkyo.ac.jp/web/geisser/p-part.pdf}{Published article}.
% Additional entries for the optional explicit dimension-two supplement.
\bibitem{Weibel}
C.~A.~Weibel, \emph{The $K$-book: An Introduction to Algebraic $K$-theory}, Graduate Studies in Mathematics \textbf{145}, American Mathematical Society, 2013. Chapter~III, Definition~5.11 and relations (D1)--(D3) fix the Dennis--Stein convention used here.
\href{https://sites.math.rutgers.edu/~weibel/Kbook/Kbook.III.pdf}{Author's Chapter~III}.

\bibitem{BunkeTamme}
U.~Bunke and G.~Tamme, \emph{Multiplicative differential algebraic $K$-theory and applications}, Annals of $K$-Theory \textbf{1} (2016), no.~3, 227--258. See Remark~2.6, Theorem~2.31, and Section~6.
\href{https://doi.org/10.2140/akt.2016.1.227}{doi:10.2140/akt.2016.1.227}.

\bibitem{Borel}
A.~Borel, \emph{Stable real cohomology of arithmetic groups}, Ann. Sci. \'Ecole Norm. Sup. (4) \textbf{7} (1974), 235--272. In particular, Proposition~12.2 gives the relevant rank calculation.
\href{https://doi.org/10.24033/asens.1269}{doi:10.24033/asens.1269}.

\bibitem{DeJeuWedge}
R.~de Jeu, \emph{Zagier's conjecture and wedge complexes in algebraic $K$-theory}, Compositio Math. \textbf{96} (1995), no.~2, 197--247. Proposition~4.1 and Remark~5.2.
\href{https://www.numdam.org/article/CM_1995__96_2_197_0.pdf}{Published article}.

\bibitem{BGTTrace}
A.~J.~Blumberg, D.~Gepner, and G.~Tabuada,
\emph{Uniqueness of the multiplicative cyclotomic trace},
Advances in Mathematics \textbf{260} (2014), 191--232, Theorem~7.4.
\href{https://arxiv.org/abs/1103.3923}{arXiv:1103.3923}.

\bibitem{BokstedtOttosen}
M.~B\"okstedt and I.~Ottosen,
\emph{A Hochschild--Kostant--Rosenberg theorem for cyclic homology},
Section~2, especially Definition~2.1, Remark~2.2, and Proposition~2.3.
\href{https://arxiv.org/abs/1601.07412}{arXiv:1601.07412}.

\bibitem{PriorIntegralK3}
OpenAI, \emph{An integral degree-three Gersten counterexample},
research preprint, 26 September 2026, p.~1.
\href{https://github.com/openai/math/blob/main/preprints/An-Integral-Degree-Three-Gersten-Counterexample-September-26-2026/paper.pdf}{Repository PDF}.

\bibitem{PriorIntegralK5}
OpenAI, \emph{An integral counterexample to Gersten's conjecture},
research preprint, 25 September 2026, p.~2.
\href{https://github.com/openai/math/blob/main/preprints/An-Integral-Counterexample-to-Gerstens-Conjecture-September-25-2026/An-Integral-Counterexample-to-Gerstens-Conjecture-September-25-2026.pdf}{Repository PDF}.

\end{thebibliography}}
\end{document}
